% 本文件由 LaTeX 排版 Agent 根据有效 translation.json 自动生成。
% author=Basil the Great; work_id=3201; title=六天

\chapter{\OriginalTerm{六天}{Hexaemeron}}

% structure: p0001 source=h1 target=omitted reason=page-title-already-rendered-by-parent

讲道一\newline{}讲道二\newline{}讲道三\newline{}讲道四\newline{}讲道五\newline{}讲道六\newline{}讲道七\newline{}讲道八\newline{}讲道九\par

\SourceRecord{Translated by Blomfield Jackson.；Nicene and Post-Nicene Fathers, Second Series；Vol. 8.；Edited by Philip Schaff and Henry Wace.；Buffalo, NY: Christian Literature Publishing Co.,；1895.；<http://www.newadvent.org/fathers/3201.htm>.}

% structure: p0001 source=h1 target=section reason=page-title

\section{\WorkTitle{创世六日（讲道篇一）}}

\BibleQuote{在起初天主创造了天地}\par

1. 任何开始叙述世界形成的人，都应当从可见事物的良好秩序开始。我即将谈论天地的创造，这并非如某些人所想象的是自发的，而是来自天主。什么样的耳朵配听到这样的故事？灵魂应当何等热切地准备自己，以接受如此崇高的教训！它应当如何纯洁，摆脱肉体的情感，不被世俗的烦扰所遮蔽，在研究上何等活跃和热切，何等渴望在周遭事物中找到关于天主的配得上祂的观念！\par

但在权衡这些说法的正当性、审视这些简短话语所含的全部意义之前，让我们看看是谁在向我们说话。因为，如果我们智慧的软弱不允许我们洞悉作者思想的深度，那么我们将因其权威的力量而不由自主地被吸引去信服他的话。现在，是梅瑟写下了这段历史；梅瑟，在襁褓中时便被形容为极其俊美；梅瑟，被法郎的女儿收养，从她那里接受王室教育，并以埃及的智者为师；梅瑟，藐视王室的荣华，宁愿与同胞一同受苦，选择同天主的百姓一同受迫害，也不愿享受罪恶的短暂欢乐；梅瑟，从本性中接受了如此对正义的热爱，以至于在领导天主百姓的重任托付给他之前，就因对邪恶的天生厌恶而驱使他追捕作恶者，甚至以死惩罚他们；梅瑟，被他曾施恩的人放逐，急忙逃离埃及的喧嚣，在埃塞俄比亚避难，远离从前的追求，在那里度过四十年，默观自然；梅瑟，最终在八十岁时看见了天主，尽人所能看见祂的程度，或者更确切地说，是先前从未赐予人看见祂的程度，按照天主亲自的见证：\BibleQuote{若你们中间有先知，我上主将在异象中向他显现，在梦中与他说话。但我的仆人梅瑟却不是这样，他在我全家是忠信的，我要与他面对面说话，明明地，而不以谜语。}就是这个被天主视为配得面对面看见祂、如同天使一般的人，将他从天主那里所学的传授给我们。那么，让我们倾听这些真理的话语，这些话语是藉着圣神的默感写成的，而非依靠\Quote{智慧动听的言辞}\BibleRef{格前 2:4}，这些话语的目的不是要博取听众的赞赏，而是要使受教者获得救恩。\par

2. \BibleQuote{在起初天主创造了天地。}\BibleRef{创 1:1}我被这个思想震撼而停住。我该先说什么？从何处开始我的故事？我该揭露外邦人的虚妄，还是高举我们信仰的真理？希腊的哲学家们费尽心思解释自然，但没有一个体系保持稳固不动摇，每个体系都被后继者推翻。驳斥他们是徒劳的；他们自己足以互相摧毁。那些因过于无知而不能认识天主的人，无法承认一个理智的起因主宰了宇宙的诞生；这一根本错误使他们陷入可悲的后果。有人求助于物质原理，将宇宙的起源归之于世界的元素。另一些人想象原子、不可分的物体、分子和管道，通过它们的结合构成可见世界的本性。原子结合或分离，产生生与死，而最持久的物体只是因相互粘附的强度而保持其一致性：这些作者编织了一张真正的蜘蛛网，他们给天、地、海如此脆弱的起源和如此不坚固的实体！这是因为他们不知道说\BibleQuote{在起初天主创造了天地。}被他们内在的无神论所欺骗，他们觉得没有什么统治或主宰宇宙，一切都听凭偶然。为防止我们犯此错误，创造论的作者从最初的话语就用天主的名字开启我们的理智：\BibleQuote{在起初天主创造了。}何等荣耀的秩序！他先确立了起初，免得人们以为世界从未有开始。然后他加上\Quote{创造}，以表明所造之物只是造物主能力的很小一部分。正如陶匠在制造了大量器皿之后，并没有耗尽他的技艺或才能，同样，宇宙的创造主，其创造能力远非局限于一个世界，而是能无限延伸，只需祂意志的推动就能使可见世界的浩瀚产生。如果世界有开始，如果它是被创造的，那么问问是谁给了它这个开始，谁又是创造者：或者更确切地说，为免人类推理使你偏离真理，梅瑟预先回答了疑问，将天主的可畏之名像印章和守护一样刻在我们心中：\BibleQuote{在起初天主创造了}——祂是慈善的本性，无限的良善，一切理性受造物值得爱的对象，最可渴望的美，万有的起源，生命之源，理智之光，不可揣测的智慧，是祂\Quote{在起初创造了天地。}\par

3. 人哪！那么，不要以为这个有形可见的世界是无始的；也不要因为天体作圆周运动，我们的感官难以确定圆周的起点，就以为受圆周运动推动的物体其本性是无始的。毫无疑问，圆周（我是指由一条线所画出的平面图形）是超出我们知觉的，我们不可能找出它从哪里起始、到哪里终结；但我们不应因此而相信它是无始的。虽然我们感觉不到，它确实从某一点开始，画图者从中心以某一半径开始画起。同样，看到作圆周运动的形体总是回归自身，一刻也不中断其规则的运行，不要徒然想象世界既无开始也无终结。因为\Quote{\BibleQuote{这世界的局面正在逝去}\BibleRef{格前 7:31}}；并且\Quote{\BibleQuote{天地要消逝。}\BibleRef{玛 24:35}}。关于终结和世界更新的教义，已经预含在这段置于受启示历史开头的简短话语中。\Quote{在起初天主创造了。}在时间内开始的事物，注定要在时间内终结。如果有开始，就不要怀疑终结。那么，几何学、算术计算、立体研究、闻名遐迩的天文学——这些劳心费神的虚空——有什么用处呢？如果追求这些学问的人以为这个有形可见的世界与万物的造物主、与天主自己同等永恒；如果他们将只属于无限不可见本性的荣耀，加给这个有物质身体的有限世界；如果他们不能理解，一个其部分有朽坏和变化的全貌，必然最终也会遭受其部分的命运？但他们却成了\Quote{\BibleQuote{他们的思念成虚妄，他们的愚昧的心就昏暗了；自称为智者，反成了愚拙。}\BibleRef{罗 1:21}}。有些人主张，天与天主自永远共存；另一些人则主张，天就是天主自己，无始无终，且是万物特殊安排的缘由。\par

4. 毫无疑问，有一天，因着这一切世俗的智慧，他们将受到更严厉的谴责，因为他们既如此清楚地看透了虚妄的学问，却故意对真理的知识闭眼不看。这些人测量星辰的距离并描述它们——无论是北方那些一直明亮照耀在我们眼中的星，还是南方居民能见到但我们却不晓得的南极星；他们将北方的星带和黄道圈分成无限多部分，精确地观察星辰的运行、它们的固定位置、它们的偏斜、它们的回归以及每颗星运行一周所需的时间；这些人，我再说，他们发现了一切，唯有一件事除外：即天主是宇宙的造物主，是按各人行为功德赏善罚恶的公正审判者。他们不懂得提升自己到万物终结的观念——那是审判教义的后果——也不懂得看到，如果灵魂从这个生命过渡到新生命，世界就必须改变。事实上，正如现世生命的本性与此世相类似，同样，在未来的生命中，我们的灵魂将享有与它们新状况相配的命运。但是他们离运用这些真理如此之远，以至于当我们向他们宣告万物的终结和时代的更新时，他们只是嘲笑。既然开始自然先于从它而来的一切，那么作者在向我们讲述那些在时间内有起源的事时，自然在叙述开头就写下这些话——\Quote{在起初天主创造了。}\par

5. 事实上，似乎即使在这个世界之前，也存在一种事物秩序，我们的心灵可以形成关于它的概念，但我们却无法言说，因为它对于一个才入门、在知识上仍是婴孩的人来说，过于高深。世界的诞生之前，有一种状态适合超自然能力的运行，超越时间的界限，永恒无限。宇宙的造物主和\OriginalTerm{工匠}{Demiurge}在其中完成了祂的工作：属神的光，为使所有爱主者幸福；理智的、不可见的本性；纯粹理智体的全部井然有序的安排——这些理智体超越我们的心灵，我们甚至无法发现它们的名字。它们充满了这个不可见世界的本质，如\Person{圣保禄}教导我们的：\Quote{\BibleQuote{一切都是藉着祂造的，在天上的和在地上的，有形的和无形的，无论是上座者、宰制者、率领者、掌权者}\BibleRef{哥 1:16}}，或是异军、天使的部队、总领天使的尊位。最后，有必要在这个世界之外再添加一个新世界，它既是灵魂受教益的学校和训练场，也是注定要出生和死亡者的家园。于是，时间相继被创造出来，其性质与这个世界以及其上生活的动植物相似，永远向前奔流、消逝，从不停止。时间难道不是这样的吗？过去的不再存在，未来的还不存在，现在的在被人认识之前就已逃逸？那生活在时间中的受造物也是如此——注定要生长或朽坏，没有安息，没有确定不移。因此，植物和动物的身体，被迫跟随某种潮流，被引向出生或死亡的运动所带走，它们应当生活在这样的环境中——其性质与那些可变化的存有相一致。因此，那位智慧地告诉我们宇宙诞生的作者，在叙述的开头就放上这些话：\Quote{在起初天主创造了；}也就是说，在时间的起初。所以，如果他在起初使世界出现，这并不证明世界的诞生先于一切被造之物。他只愿告诉我们：在不可见的、理智的世界之后，有形可见的、感官的世界开始存在。\par

第一个运动被称为\Quote{开始}。\Quote{行正义是善道的开始。}正义的行为确实是迈向幸福生活的第一步。同样，我们也把事物的本质和最初部分称为\Quote{开始}，例如房屋的地基、船只的龙骨；正是在这个意义上，经上说：\BibleQuote{敬畏主是智慧的开始，} \BibleRef{箴 9:10} 也就是说，虔敬仿佛是成全的根基和基础。技艺也是艺术家作品的开端，贝匝肋耳的技巧\Emph{开始}了会幕的装饰。甚至作为终极原因的美善也常常是行为的\Emph{开端}。因此，天主的赞许是施舍的开端，以及那为我们存留在恩许中的终向是一切美德努力的开端。\par

6. \Quote{开始}一词既然有这些不同的含义，请看看我们在此是否包含了所有这些意义。你可以知道这个世界开始形成的时间：如果你追溯过去，努力发现第一天，你就会找到时间的第一个运动；然后，天地受造仿佛根基和基础；接着，如\Quote{开始}一词所指出的，一种理智的理性在可见事物的秩序中主持。最终你会发现，世界不是偶然地、无理性地被设想出来的，而是为了一个有益的目的和一切受造物的重大利益；因为它实在是理性灵魂的学校，在其中操练，训练场，在那里学习认识天主；因为通过视见可见和可感觉的事物，心智如同被手牵引，达到对不可见之物的默观。\Quote{因为，}如宗徒说：\BibleQuote{自从受造以来，他那看不见的美善，都可凭他所造的万物，辨识洞察出来。} \BibleRef{罗 1:20} 也许\BibleQuote{在起初天主创造了}这句话，意味着创造的迅速且不可察觉的瞬间。实际上，开始是不可分割且瞬间的。路的开端还不是路，房屋的开端还不是房屋；同样，时间的开端还不是时间，甚至连它的最小部分也不是。如果某个反对者告诉我们，开始是一个时间，那么他显然应当将其分成时间的各部分——开始、中间和结束。而想象一个开始的开始是荒谬的。此外，如果我们把开始分成两个，我们就得到两个而不是一个，或者毋宁说得到多个，实际上得到无限多，因为一切可分割的东西都可以无限分割。所以，如果经上说：\BibleQuote{在起初天主创造了}，这是要教导我们，在天主的意愿下，世界在不到一刹那的时间内兴起；其他译者更清楚地表达了这一含义，他们说：\Quote{天主概略地造了}，也就是说，立刻在一瞬间。关于开始，从众多要点中只提出几点，就足以说明问题了。\par

7. 在艺术中，有些旨在生产，有些旨在实践，有些旨在理论。后者的对象是思想的操练，第二类的对象是身体的运动。如果停止，一切便静止；再也看不到什么。因此，舞蹈和音乐没有后续；它们没有自身以外的目的。相反，在创造性艺术中，作品在操作之后仍然存在。例如建筑——以及那些在木、铜和纺织中工作的艺术，实际上，所有那些即使在工匠消失后，仍显示其勤勉的智慧，并使建筑师、铜匠或织工因其作品而被赞赏的艺术。因此，为了表明世界是一件为众人观赏而展示的艺术品，为了让他们认识创造它的那一位，梅瑟没有使用别的词。他说：\BibleQuote{在起初天主创造了。}他没有说\Quote{天主工作了}、\Quote{天主形成了}，而是说\Quote{天主创造了。}在那些想象世界从永远就与天主共存的人中，许多人否认世界是由天主创造的，而说它自发地存在，如同这能力的阴影。他们说，天主是其因，但却是非自愿的因，如同身体是阴影的因，火焰是光亮的因。正是为了纠正这个错误，先知如此精确地陈述：\BibleQuote{在起初天主创造了。}他没有使事物本身成为其存在的原因。祂是善的，便造了有用的作品；祂是智慧的，便造了最美的一切；祂是强有力的，便造了极其伟大的。梅瑟几乎向我们展示了至高工匠的手指，握住了宇宙的实体，使不同的部分完美和谐地形成，并使整体产生出一曲和谐的乐章。\par

\BibleQuote{在起初天主创造了天地。}通过命名这两个极端，他暗示了整个世界的实质，给了天上首位的特权，把地放在第二等级。所有中间的存在物都与极端同时受造。因此，虽然没有提到元素——火、水、气——但想象它们都混合在一起，你会发现水、气和火都在地上。因为火从石头中迸出；从地里挖出的铁在摩擦下产生大量的火。一个奇妙的事实！火封闭在物体中，潜伏在那里而不伤害它们，但一旦释放，它就吞噬那先前保存它的东西。地含有水，正如挖井者教导我们的。它也含有气，正如在太阳的热力下，潮湿时地上散发出的蒸汽所显示的。既然按其本性，天在空间中占据高处，地占据低处（实际上，我们看到一切轻的东西向天上升，重的东西落向地面）；既然高度和深度是最彼此对立的点，那么提到最远的部分就足以表示所有填充中间空间的事物的包含。因此，不要要求一一列举所有元素；从圣经的提示中，猜测一切被默默略过的事物。\par

8. \Quote{在起初天主创造了天地。} 如果我们想要发现每一个呈现在我们默观中或落入我们感官之下的存有的本质，我们就会陷入冗长的离题话，而彻底审视这个问题所需的言辞将超出我的能力。在这些问题上花费时间不会证明对教会有益。关于天的本质，我们满足于依撒意亚所说的，因为他用简单的语言给了我们对其本性足够的认识：\Quote{天如烟云造成}，也就是说，祂创造了一种微妙的、没有固体性或密度的物质，用以形成天。至于天的形状，我们也满足于同一位先知的话，当他赞美天主\Quote{铺张诸天如幔子，展开诸天如可居住的帐棚}。同样，关于地，让我们下定决心不要因试图探求其本质而折磨自己，不要因寻找它所隐藏的实体而耗尽我们的理性。不要让我们寻求任何按其存在条件而缺乏属性的本性，但让我们知道，我们看见它所披戴的一切现象都与其存在的条件相关并完善其本质。试着用理性拿走它所拥有的每一种属性，你将一无所获。拿掉黑色、冷、重量、密度、与味觉有关的属性，总之是我们看见它所有的这一切，实质就消失了。\par

如果我请你放下这些空洞的问题，我不会期望你试图找出地的支撑点。理性会因看到自己的推论迷失在无限中而感到眩晕。你说地安卧在一层空气上？那么，这柔软且无实体的物质如何能承受压在其上的巨大重量？它如何能不向各方滑移，以避免下沉的重量，并覆盖在压迫它的质量之上？你假定水是地的根基？那么你将总是不得不问自己，如此沉重且不透明的物体如何不穿过水；如此庞大的质量如何被比自身更弱的本性所支撑。然后你必须为水寻找一个基础，并且你会很难说出水本身安放在什么之上。\par

9. 你假定一个更重的物体阻止地掉入深渊？那么你必须考虑这个支撑本身需要一个支撑以防止它掉落。我们能想象一个吗？我们的理性再次要求另一个支撑，如此我们就会陷入无限，总是为已经找到的基础再想象一个基础。并且我们越深入这个推理，就越不得不赋予这个基础更大的力量，以便它能支撑所有压在其上的质量。那么为你的思想设定一个界限，以免你探究不可理解之事的好奇心招致约伯的责备，并且不被祂问：\Quote{地的根基安置在何处？} \BibleRef{约 38:6} 如果你曾在圣咏中听到\Quote{我使地的柱子稳固}，请在这些柱子中看见支撑它的力量。因为另一段经文\Quote{祂将地奠定在海上}如果不是指水环绕着地四周，又是什么意思呢？那么，水这流动的元素，沿着每一处斜坡流下，如何能悬停而不流动呢？你没有想到，地自行悬置的想法会使你的理性陷入类似但更大的困难，因为从其本性而言它更重。但让我们承认地自托于自身，或者说它乘驾在水上，我们仍必须忠于真实敬虔的思想，并认识到一切都由造物主的能力所支撑。那么让我们回答自己，也回答那些问我们这巨大质量靠什么支撑的人：\Quote{地的四极在祂手中。} 这是一个既为我们自己的教导，也为听众有益的绝对无误的教义。\par

10. 有一些自然的探究者，用大量华丽的言辞给出地不动的理由。他们说，地置于宇宙的中心，且不能向一边倾斜多于另一边，因为它的中心处处与表面等距，所以它必然自托于自身；因为各处相等的重量不能倾向任何一边。他们接着说，地占据宇宙的中心并非无因或偶然，而是其自然且必然的位置。正如天体占据空间的最高端点，所有重物，他们论证说，我们可以假设从这些高区域掉落的，都会从各个方向被带到中心，而各部分趋向的点显然就是整个质量将被推聚的地方。如果石头、木头、所有地上的物体都从上往下落，那么这必定是整块地的恰当且自然的位置。相反，如果一个轻体离开中心，显然它将向更高区域上升。因此，重物从顶部向底部运动，而根据这个推理，底部正是世界的中心。那么，不要惊讶世界永不坠落：它占据宇宙的中心，其自然位置。它必然被迫留在原位，除非一个违背自然的运动使其移动。如果在这个体系中有什么对你显得可能，请为如此完美秩序的源头，为天主的智慧保持你的钦佩。宏伟的现象不会因为我们对它们奇妙机制有所发现而减少对我们的震撼。这里难道不同吗？无论如何，让我们更喜爱信仰的单纯而非理性的论证。\par

11. 我们也可对诸天说同样的话。此世的智者以何等的喧哗议论过它们的本质！有人说，天是由四种元素组成的，因为它们是有形可见的，并且因抵抗力而有土、因目击而有火、因混合而有气与水。另一些人则因这体系似不可信而拒绝之，并按其自己的构想为世界引入了第五种元素来形成诸天。他们说，存在一种以太体，它不是火、气、土、水，简言之也不是任何单纯物体。这些单纯物体有其自然的直线运动：轻者向上，重者向下；而此种上下运动与圆周运动并不相同，直线与圆周之间有着最大的差异。因此，运动如此不同的物体必然在本质上也有差异。但是，甚至不可能假定诸天是由我们称为元素的原始物体构成的，因为相反力量的结合无法产生均匀自发的运动，因为每个单纯物体从自然获得不同的冲动。因此，使复合物体保持连续运动是费力的，因为不可能使它们任何一个运动与所有那些不和谐的运动协调一致；因为轻质粒子的本性是与较重粒子的本性相争的。若我们试图上升，就会被土元素的重量阻止；若我们向下投掷，就会强迫我们中的火性部分违背其本性往下拖。这种元素的斗争导致了它们的解体。一个受到暴力且被置于与自然对立位置的物体，在经过短暂但有力的抵抗后，会很快分解成与它原有元素同样多的部分，每个组成部分返回其自然位置。正是这些理由的力量（诸天和星辰之第五种物体的发明者说）迫使他们拒绝前人的体系，而诉诸他们自己的假说。然而，又有一位雄辩家起来，驱散并摧毁了这一理论，以便让自己发明的观念占优势。\par

不要让我们追随他们，以免陷入同样的轻浮；让他们彼此驳斥吧，而我们，不必为本质而烦恼，应与梅瑟同声说：\BibleQuote{天主创造了诸天和大地。}让我们为一切明智而巧妙制成之物赞颂至高艺匠；藉着可见之物的美丽，让我们提升自己到超越一切美丽的那一位；藉着有形且有限之身体的宏大，让我们思议无限的存有，祂的无限和全能超越一切想象的效力。因为，尽管我们忽略受造物的本质，但那些从各方向吸引我们注意的对象是如此奇妙，以致最敏锐的理智也无法达到认识世上最小现象的知识，无论是为了给予恰当的解释，还是为了向造物主献上应得的赞美，祂拥有一切荣耀、一切尊威和一切权能，永世无尽。阿们。\par

\SourceRecord{Translated by Blomfield Jackson.；Nicene and Post-Nicene Fathers, Second Series；Vol. 8.；Edited by Philip Schaff and Henry Wace.；Buffalo, NY: Christian Literature Publishing Co.,；1895.；<http://www.newadvent.org/fathers/32011.htm>.}

% structure: p0001 source=h1 target=section reason=page-title

\section{\WorkTitle{Hexaemeron（第二篇讲道）}}

\Emph{\Quote{\Emph{地是看不见的，且未成形。}}}\par

1. 在今日早晨占据我们的寥寥数语中，我们发现了如此深奥的思想，以至于我们绝望于进一步深入。如果这只是圣所的前院，如果圣殿的门廊如此宏伟壮丽，如果其美丽的光彩如此炫目灵魂之眼，那么至圣所又将如何？谁敢试图进入最内层的圣所？谁能窥探其中的奥秘？确实，凝视它是被禁止的，而语言也无能为力来表达心灵所领悟的。然而，既然公义的审判者为单单行善的意向存留了赏报，且是最可欲的赏报，让我们不要迟疑继续我们的探究。即使我们可能无法达到真理，但若藉着圣神的帮助，我们不离弃圣经的意义，我们就不该被拒绝，并且藉着恩宠的帮助，我们将有助于天主的教会的建设。\par

\BibleQuote{地}，圣经说，\BibleQuote{是看不见的，且未成形。}天地被造时并无区别。那么，为何天是完美的，而地仍是未成形的和不完备的？一句话，地的不完备状态是什么？又为何它是看不见的？地的肥沃就是它的完美完成；各种植物的生长，高大树木的萌发，无论是结果的还是不结果的，花朵的甜香和美丽的色彩，以及稍后，在天主的声音下，从地中出来以美化她的万物，她们共同的母亲。由于这一切尚不存在，圣经称地为\BibleQuote{无形}是正确的。我们也可以说，那时天还未完成，尚未得到它天然的装饰，因为它尚未被太阳和月亮的光辉照耀，也未被星辰的合唱加冕——这些天体尚未被造。因此，你称天也是\BibleQuote{无形}，并不偏离真理。地看不见有两个原因：或许是因为观看者——人——尚未存在，或者因为地沉浸于覆盖其表面的水下而无法被看见，因为水尚未被聚集到它们自己的地方——后来天主将它们收集在那里，并称之为海。什么是看不见的？首先，我们肉眼不能感知的一切，例如我们的心灵；其次，本性可见却被某些物体遮蔽的东西，如深藏地下的铁。正是在这个意义上，因为它被水隐藏，地仍是看不见的。然而，既然光尚未存在，且地处于黑暗之中，因为其上空气的幽暗，我们不应惊讶圣经因这个原因称它为\BibleQuote{看不见的}。\par

2. 但真理的败坏者，不能使他们的理智服从圣经，随意歪曲圣经的意义，他们声称这些词指物质。因为他们说，物质就其本性而言是无形的和看不见的——按其存在条件，没有性质，没有形式和形状。工匠藉着将其置于祂智慧的运作之下，给它披上形式，组织它，从而赋予可见世界以存在。\par

如果物质是非受造的，它就有权要求与天主同等的尊荣，因为它必须与祂同等。这不是邪恶的顶峰吗？一种极端的畸形，没有性质，没有形式、形状，没有构形的丑陋——用他们自己的说法——竟然与祂——祂是智慧、能力和美本身，宇宙的创造者和\OriginalTerm{造物主}{\GreekText{δημιουργός}}——享有同等的特权？这还不是全部。如果物质如此伟大，足以被天主的全部智慧所作用，它就会在某种程度上将其\OriginalTerm{位格}{\GreekText{ὑπόστασις}}提升到与天主不可及的能力同等，因为它能用自身衡量神圣理智的全部范围。如果它对天主的运作来说是不足的，那么我们就陷入了更荒谬的亵渎，因为我们谴责天主，不能因物质的缺乏而完成祂自己的工程。人类本性的贫乏欺骗了这些推理者。我们的每种技艺都作用于某种特殊的材料——铁匠的技艺作用于铁，木匠的技艺作用于木。在一切之中，都有主体、形式和由形式产生的工作。物质是从外部取来的——技艺赋予形式——而工作同时由形式和物质构成。\par

这就是他们为自己设想的天主化工的概念。世界的形态源于至高工匠的智慧，而物质则从外部来到造物主面前；因此，世界产生于双重起源：它从外部获得了物质和本质，从天主获得了形态和样貌。他们由此否定全能的天主掌管了宇宙的成形，却声称祂只对一个共同的作品做出了关键的贡献，祂只对万有的生成贡献了一小部分；他们因自己推理的卑劣而无法将目光提升至真理的高度。在此世，技艺是后于物质的——因不可或缺的需要而被引入生活。羊毛在纺织用它弥补自然的一个缺陷之前就已存在。木材在木工掌握它并每日改造它以供应新需求之前就已存在，使我们看到由此产生的种种益处：为水手提供船桨，为劳动者提供扬谷扇，为士兵提供长矛。但是，在那些如今吸引我们注意的一切事物存在之前，天主在祂心中权衡并决定使尚无时间的事物进入存在时，就按世界应有的样子想象了它，并按祂希望赋予它的形态创造了物质。祂为诸天指定了适合诸天的本性，为大地赋予了与其形态相应的本质。祂按自己的意愿形成了火、空气和水，并为每一种赋予了其存在目的所需的本质。最后，祂以不可分解的联结锻造了宇宙的一切不同部分，并在它们之间建立了如此完美的共融与和谐，以至最遥远的距离，尽管相距遥远，也显得在一种普遍的同情中联合在一起。因此，让那些人放弃他们神话般的想象，他们尽管论证薄弱，却妄图衡量一种在人的理性中不可理解、在人的言语中不可言说的能力。\par

3. 天主创造了诸天和大地，但并非只创造一半——祂创造了整个诸天和整个大地，即连同形态一起创造了本质。因为祂不是形态的发明者，而是万有本质的创造者。进一步，让他们告诉我们，天主的能动能力如何处理物质被动的本性：后者提供无形态的物质，前者拥有无物质的形态知识，两者彼此需要——创造者为了展示祂的技艺，物质为了不再无形态并接受形态。但让我们在此停止，回到我们的主题。\par

\Quote{\Emph{大地是无形且未完成的。}}圣作者在说\BibleQuote{在起初天主创造了诸天和大地}时，略过了许多东西：水、空气、火以及它们的结果——这些实际上构成了世界的真正完善，无疑是与宇宙同时被造的。通过这种沉默，圣经历史愿训练我们理智的活动，给它一个微弱的起点，以促其发现真理。因此，我们没有被告知水的创造；但是，既然我们被告知大地不可见，请问是什么可能覆盖了它，阻止它被看见？火不能隐藏它。火照亮一切，散播光明而非黑暗。空气也不能包裹大地；空气本性稀薄透明，它接收所有可见物体并将其传递给观看者。只剩下一种假设：漂浮在大地表面的是水——尚未被限制在自身之处的流动本质。因此，大地不仅不可见，而且尚未完成。即使今天，过多的湿气也阻碍大地的丰产。同一原因同时阻止了它被看见和完成，因为大地恰当而天然的装饰就是其完成：山谷中摇曳的谷穗——长满青草、缀以多彩花朵的草地——由森林遮荫的肥沃林间空地和小丘。这一切尚未产生；大地正凭借从造物主那里获得的能力孕育着它们，但它在等待指定的时间和神圣的命令来分娩。\par

4. \BibleQuote{黑暗笼罩深渊的面。}\BibleRef{创 1:2} 如果人任由自己的幻想扭曲这些话语的意义，那便是新的寓言和极不敬的想象之来源。这些邪恶之人不按实际意义——未被照亮之空气、物体遮挡产生的阴影，或因某种原因缺乏光明的地方——来理解\Quote{黑暗}。对他们而言，\Quote{黑暗}是一种邪恶的力量，或者更确切地说，是邪恶的化身，它在自身内起源，与天主的良善对立并永恒斗争。他们说，如果天主是光，那么毫无疑问与之斗争的力量必是黑暗；\Quote{黑暗}并非来自外在根源，而是自行存在的邪恶。\Quote{黑暗}是灵魂的仇敌，死亡的主要原因，德行的对手。他们在谬误中说，先知的话语表明它存在，且并非来自天主。由此幻想出了何等乖谬邪恶的信理！何等残暴的豺狼\BibleRef{宗 20:29}撕裂主的羊群，从这些话语中跳出扑向灵魂！难道不是从这产生了马西翁们和瓦伦提努们，以及可憎的摩尼教异端——你若不甚谬误，可称之为教会的腐败体液？\par

人啊，你为何如此偏离真理，为自己想象那将导致你灭亡的事？这话简单明了，人人都能领会。\Quote{大地是看不见的。}为何？因为\Quote{深渊}覆盖了它的表面。什么是\Quote{深渊}？是极深的水体。但我们知道，透过清澈透明的水，我们能看见许多物体。那么，为何没有一部分大地透过水显现出来？因为环绕它的空气仍没有光，处于黑暗之中。太阳的光线穿透水，常常让我们看见河床的卵石，但在漆黑的夜里，我们的视线无法穿透水面。因此，\Quote{大地是看不见的}这句话可由后面的话来解释：\Quote{深渊}覆盖了它，而深渊本身在黑暗中。所以，深渊并非如人所想象的那样是一群敌对势力，\Quote{黑暗}也不是与善为敌的邪恶主权力量。事实上，两种力量相等的敌对原则，若不断相互攻击，最终会自我毁灭。但如果一方取得优势，就会完全消灭被征服者。因此，在善恶斗争中说它们保持平衡，就等于说它们在进行一场无休止的战争和永恒的毁灭，而对立双方同时既是征服者又是被征服者。如果善更强，有什么能阻止恶被完全消灭？但若如此（说这话已是亵渎），我问自己，他们自己怎能不充满恐惧，竟想象出如此可憎的亵渎之词？\par

同样，说恶起源于天主也是亵渎；因为对立者不能从其对立方产生。生命不产生死亡；黑暗不是光的起源；疾病不是健康的制造者。在状态的改变中，有一种状态过渡到相反的状态；但在生成中，每个存在者都从与其相似者而来，而非从其对立者而来。那么，如果恶既不是非受造的，也不是由天主创造的，它的本性从何而来？诚然，恶的存在，世上没有一个人会否认。那我们将怎么说呢？恶不是一种有生命的、有灵性的本质；它是灵魂违背德行的状态，因疏忽而远离善所产生的。\par

5. 那么，不要超越自身去寻找恶，也不要想象有一种原始的邪恶本性。我们每个人都应当承认，自己是自身恶习的首要作者。在日常生活中的事件中，有些是自然发生的，如衰老和疾病；有些是偶然的，如意外事件，其原因超出我们自身，常常令人悲伤，有时幸运，例如挖井时发现宝藏，或去市场时遇到疯狗。另一些则取决于我们自己，如管束自己的激情，或不对自己的快乐加以约束，控制自己的愤怒，或向激怒我们的人动手，说实话或说谎，拥有温和有节的性情，或暴躁、傲慢、膨胀。在这里，你是自己行为的主人。不要在自己之外寻找主导原因，而要承认，恶（严格来说）除了我们自愿的堕落之外，没有其他起源。如果它是不自愿的，不取决于我们自己，法律就不会对罪人如此恐惧，法庭在根据罪行轻重判决可怜人时，也不会如此毫不留情。但关于严格意义上的恶，说得够了。疾病、贫困、卑微、死亡，以及人类的一切痛苦，都不应被视为恶；因为我们不把它们的对立面视为最大的恩赐。在这些痛苦中，有些是自然的结果，另一些显然对许多人有益。那么，我们现在暂且不谈这些比喻和寓意，而是简单地、不带无谓好奇地遵循\WorkTitle{圣经}的话，从黑暗本身去理解它给我们的启示。\par

但理性会问：黑暗是与世界同时被创造的吗？它比光更古老吗？为什么尽管它低劣，却先于光存在？我们回答说：黑暗并非本质的存在，而是光离去时在空气中产生的一种状态。那么，从世界消失的光是什么，使得黑暗覆盖了深渊的表面？如果在可感觉的、可朽坏的世界形成之前有任何东西存在，我们无疑会断定它存在于光明中。天使的品级、天上的万军、所有已知或未知的理智本性、所有服役的神灵，并非生活在黑暗中，而是享有适合他们的光明和属灵喜乐的状态。\par

没有人会否认这一点；尤其是那些盼望天上之光作为许诺给德行的赏报之一的人，那光正如\Person{撒罗满}所说，永远是义人的光，那光使宗徒说：\Quote{感谢圣父，祂使我们有资格在光明中分享圣人的产业。}\BibleRef{哥 1:12}最后，如果被定罪的人被投入外面的黑暗，显然那些蒙天主悦纳的人安息于天上的光明。那么，按照天主的命令，天穹出现，包裹其圆周内的一切，一个巨大而不可分割的物体将外部与内部隔开，它必然使内部空间因缺乏与外光的沟通而陷入黑暗。的确，形成阴影需要三样东西：光、物体和暗处。天穹的阴影形成了世界的黑暗。我请你通过一个简单的例子来理解我的意思：正午时分，你搭起一个由致密不透光的材料制成的帐篷，将自己关在里面，就会突然陷入黑暗。假设最初的黑暗就是这样，并非直接自存，而是由某些外在原因造成。若说它停留在深渊之上，是因为空气的末端自然接触物体的表面；而当时水覆盖一切，我们就不得不说黑暗在深渊的面上。\par

6. \Emph{\BibleQuote{天主之神运行在水面上}}。这\Quote{圣神}是指空气的弥漫吗？圣经作者想要向你列举世界的元素，告诉你天主创造了天、地、水、空气，而空气当时已弥漫并运动；或者更真实且经古人权威肯定的是，藉天主之神，他意指圣神。正如已指出的，这是圣经喜悦给予圣神的特殊名称，超乎其他一切名称，而\Quote{天主之神}总是指圣神，即那圆满天主圣三的圣神。因此，你最好按此意义理解。那么，天主之神如何运行在水面上呢？我即将给你的解释并非我的原创，而是一位叙利亚人的解释，他在此世的智慧上无知，却在真理的知识上精通。他说，叙利亚文词语更具表现力，且更接近希伯来文词语，从而更贴近经文意义。该词的意义是：叙利亚人认为，\Quote{运行}意味着它滋养了水的本性，如同人们看见鸟用身体覆盖蛋并从自身温暖给予它们活力。这正是这些词尽可能接近的意思——\Quote{神运行}：让我们理解，即预备了水的本性以产生生物；对于那些问圣神是否积极参与了世界创造的人，这是一个充分的证明。\par

7. \Emph{\BibleQuote{天主说：要有光}}。\BibleRef{创 1:3} 天主的第一个言语创造了光的本性；它使黑暗消散，驱除阴霾，照亮世界，同时赋予万有甜美而优雅的面貌。直到那时还笼罩在黑暗中的天空，展现出我们今天仍能看到的美。空气被点亮，或者不如说使光与它的实体混合循环，迅速向各方散布其光辉，直至极端边界。它向上直冲以太和苍穹。瞬间照亮了世界的整个幅度：东、西、南、北。因为以太也是一种如此精微透明的实体，以至于光穿过它不需要一瞬间。正如它使我们的视线瞬间到达所见之物，同样，在毫无间隔、以思想无法企及的速度下，它接收光线至其最远边界。有了光，以太变得更悦目，水变得更清澈。后者不仅接受其光辉，还通过光的反射将闪烁的光芒向各方送出。天主的言语使每样物体更愉悦、更吸引人，如同深海中的人倾注油，使他们周围的地方明亮。因此，以一个言语、在一瞬间，万有的创造者将光的恩赐赐予世界。\par

\Emph{\BibleQuote{要有光}}。这命令本身就是一种运作，一种状态被带入了存在，对此人的心智甚至无法想象更能让我们享受的。必须充分理解，当我们谈论天主的声、言语、命令时，这种神性语言并不指从发声器官逸出的声音、空气被舌头撞击；它是天主意志的简单标志，而如果我们赋予它命令的形式，只是为了更好地打动我们教导的灵魂。\par

\Emph{\BibleQuote{天主看了光，认为好}}。\BibleRef{创 1:4} 在造物主为光作的见证之后，我们怎能恰当地赞美光呢？即使在人类中，言语也将判断交于眼睛，无法提升到感官已接受的概念。但如果物体之美在于部分的比例和颜色的和谐，那么在像光这样简单而同质的本质中，如何保持美的概念？光的对称性难道不更少体现在其部分，而更多体现在看到它时的愉悦和喜悦吗？金子也是如此，它的美不是由于各部分幸运的混合，而只是因其美丽的颜色具有吸引眼睛的魅力。\par

同样，金星是最美的星辰：并非因为它的各部分组成和谐整体，而是由于它纯净而美丽的亮度映入眼帘。此外，当天主宣告光的\Quote{好}时，并非着眼于眼睛的魅力，而是为了未来的益处，因为那时还没有眼睛能判断它的美。\Quote{\Emph{\BibleQuote{天主将光与黑暗分开}}}；\BibleRef{创 1:4} 也就是说，天主赋予它们不能混合、永恒对立的本性，并在它们之间放置了最广阔的空间和距离。\par

8. \Quote{\Emph{\BibleQuote{天主称光为\InnerQuote{昼}，称黑暗为\InnerQuote{夜}}}}。\BibleRef{创 1:5} 自从太阳诞生以来，它散布在空气中的光当照耀我们的半球时为昼，而因它消失产生的阴影为夜。但在那时，不是根据太阳的运行，而是按照这种原始的光散布在空气中或按天主规定的度量收回，白昼来临，继而是黑夜。\par

\BibleQuote{\Emph{晚上过了，早晨来了，这是第一天。}}\BibleRef{创 1:5} 因此，晚上是白昼与黑夜的共同边界；同样，早晨构成黑夜走向白昼的过渡。圣经为了赋予白昼长子的特权，将第一天的结束放在第一夜之前，因为黑夜跟随白昼：在创造光之前，世界并非在黑夜中，而是在黑暗中。它是白昼的对立面，被称为黑夜，并且直到白昼之后才获得它的名称。晚上和早晨就这样被创造了。圣经所指的是一天一夜的间隔，此后不再说起白昼和黑夜，而是以更重要者的名称来称呼两者：这是一个你将贯穿圣经始终发现的惯例。处处都以天数计算时间，而不提黑夜。圣咏作者说：\Quote{我们的岁数如同日子。} 雅各伯说：\BibleQuote{我生活的年岁又少又苦，}\BibleRef{创 47:9} 其他地方还说\Quote{我一生的时日。} 因此，在历史的形式之下，为后来之事制定了律法。晚上和早晨是一天。为什么圣经说\Quote{一天是第一日}？在向我们说起第二日、第三日、第四日之前，称那开始序列的一天为\Quote{第一}难道不是更自然吗？因此，如果说它称\Quote{一天}，那是为了确定昼与夜的度量，并将它们所包含的时间结合起来。二十四小时填满了一天的间隔——我们指的是一天一夜；如果在二至点时它们并不等长，但圣经所指的时间并没有限制它们的长度。它仿佛在说：二十四小时丈量了一天的间隔，或者实际上，一天是太阳从一点出发再返回那一点所需的时间。因此，每当在太阳的运转中，晚上和早晨占据世界时，它们的周期性交替从未超过一天的间隔。但是我们必须相信这有一个奥妙的理由吗？天主创造了时间的本性，以日的间隔来测量和确定时间；并且，愿意赋予一周作为度量，祂命令一周从时期到时期自行旋转，以计算时间的运动，形成由一日旋转七次而成的一周：一个适当的圆开始并结束于自身。永恒也是如此，在自身内旋转，不结束于任何地方。因此，如果时间的开始被称为\Quote{一天}而非\Quote{第一日}，那是因为圣经愿意确立它与永恒的关系。实际上，称那本性完全分离并与其他日子孤立的日子为\Quote{一天}是恰当和自然的。如果圣经向我们说起许多世代，处处说\Quote{世代世代，万世万代}，我们并没有看到它枚举它们为第一、第二、第三。由此我们看出，这里显示的与其说是世代的界限、终结和更迭，不如说是不同状态和行动方式的区别。圣经说：\BibleQuote{上主的日子是伟大而极其恐怖的，}\BibleRef{岳 2:11} 其他地方说：\BibleQuote{祸哉，那些盼望上主日子的人！上主的日子对你们有什么好处？上主的日子是黑暗而不是光明。}\BibleRef{亚 5:18} 对于那些配得黑暗的人来说，那日子是黑暗的。不；这个没有晚上、没有更迭、没有终结的日子在圣经中并非未知，这就是圣咏作者所称的第八天，因为它不在这个以周计算的时间之外。因此，无论你称它为日子，还是永恒，你所表达的是同一个观念。给这种状态冠以日子的名称；并没有多个，只有一个。如果你称之为永恒，它仍然是独一无二的，而不是多重的。因此，为了引导你的思想朝向未来的生命，圣经用\Quote{一天}这个词语来标记那作为永恒预像的日子，日子的初果，与光同时的，因我们的主的复活而受尊崇的圣主的日子。\Emph{晚上过了，早晨来了，这是一天。}\par

但是，当我与你们谈论世界的第一晚上时，晚上出其不意地降临，结束了我的讲论。愿真光的父，祂以天光装饰了白昼，祂使火焰在夜间照耀我们，祂在来世的平安中为我们保留着属灵的永恒之光，光照你们的心灵，使你们认识真理，保护你们不致跌倒，并恩赐你们\BibleQuote{在白天端正行事}\BibleRef{罗 13:13}。这样，你们将如同太阳在诸圣的光荣中闪耀，而我将因你们在基督的日子内夸耀，愿一切荣耀和权能归于祂，直到永远。阿们。\par

\SourceRecord{Translated by Blomfield Jackson.；Nicene and Post-Nicene Fathers, Second Series；Vol. 8.；Edited by Philip Schaff and Henry Wace.；Buffalo, NY: Christian Literature Publishing Co.,；1895.；<http://www.newadvent.org/fathers/32012.htm>.}

% structure: p0001 source=h1 target=section reason=page-title

\section{\WorkTitle{六日创造（第三篇讲道）}}

\Emph{论苍穹。}\par

1. 我们现已讲述了第一天的工程，或者更确切地说，一天之内的工程。我决不敢剥夺它享有的特权，即成为创造主专属的一天，一天不与其余日子同列的日子。我们昨天的讨论论述了这天的工程，并将叙述分成两部分，早晨给你们灵魂的食粮，晚上给你们喜乐。今天我们转向第二天奇迹的默想。在此，我不愿谈论叙述者的才能，而是谈论圣经的恩宠，因为叙述如此自然，令人喜悦并得以安慰一切真理的朋友。圣咏作者如此有力地表达真理的这种魅力，他说：\BibleQuote{你的言语在我上颚何等甘美，在我口中比蜜更甜。} 昨天，我们尽己所能，通过谈论天主的圣言滋养了灵魂，今天第二天我们再次聚集，默想第二天的奇迹。\par

我知道许多从事机械行业的工匠正簇拥在我周围。一天的劳动几乎不足以维持他们生计；因此我不得不缩短我的讲论，以免使他们长时间离开工作。我该对他们说什么？你们借给天主的时间不会丢失：祂必会连本带利地归还你们。无论什么困难困扰你们，主必会驱散它们。对于那些选择神性福祉的人，祂将赐予身体的健康、心智的敏锐、事业的成功和连绵的繁荣。而且，即使在此生我们的努力未能实现我们的希望，圣神的教导对于将来的世代仍然是一笔丰富的宝藏。因此，将你们的心从今生的挂虑中释放出来，专心致志于我的话。若你们身体在此，心却挂虑地上的财富，于你们何益？\par

2. 天主说：\BibleQuote{在水与水之间要有苍穹，将水分开。} \BibleRef{创 1:6} 昨天我们听到天主的谕令：\BibleQuote{要有光。} 今天则是：\BibleQuote{要有苍穹。} 似乎这其中还有更多含义。这话并不限于一个简单的命令。它道出了苍穹结构所必需的理由：它说，是为了分开水与水之间。首先，让我们问：天主如何说话？是按我们的方式吗？祂的理智从对象接受印象，然后构想它们，并用每种对象特有的符号表达出来？祂因此必须借助发声器官来传递祂的思想吗？祂必须用声音的清晰动作敲击空气，以揭示隐藏在祂心里的思想吗？说天主需要如此曲折的方法来表达祂的思想，岂不像一个无聊的寓言？更符合真正敬礼的说法难道不是：天主的旨意和神圣理智的最初推动就是天主的圣言？圣经模糊地描述了祂，向我们表明天主不仅愿意创造世界，而且愿意借助一位同工者来创造。圣经本可以按开始的方式继续历史：起初，天主创造了天和地；此后，祂创造了光；然后，祂创造了苍穹。但是，通过让天主命令和说话，圣经含蓄地向我们显示了接受这些命令和言语的那一位。这并不是因为它吝啬于让我们认识真理，而是为了通过向我们显示一些奥秘的痕迹和暗示来点燃我们的渴望。我们欣喜地抓住并小心保存通过艰辛努力获得的果实，而轻易获得的拥有则被轻视。这就是圣经引导我们认识独生子的道路和路径。当然，天主无形体的本性不需要声音的物质语言，因为祂的思想可以直接传递给祂的同工者。那么，对于那些只凭思想就能彼此沟通意愿者，语言有何必要？声音是为听觉而造，听觉是为声音。在没有空气、舌头、耳朵以及通向头部感觉中枢的曲折通道的地方，语言是不需要的：灵魂的思想足以传递意愿。正如我所说，这种语言只是一种明智而巧妙的构思，以激发我们寻找话语所指向的那一位。\par

3. 其次，那被称为天的穹苍与起初天主所造的穹苍是否有别？有两个天吗？那些讨论天的哲学家们宁愿失去舌头也不肯承认这一点。他们声称只有一个天，其本性既不容许有第二个，也不容许有第三个，更不容许有许多个。他们说，天体的本质完全构成了其广大的统一体。因为他们说，每一个圆周运动的天体都是单一而有限的。如果这个天体被用于建造第一个天，那么就没有剩余的材料用来创造第二个或第三个天。这里我们看到了那些将未受造的物质置于造物主手下的人的想象；这是一个源自第一个寓言的谎言。但我们请求希腊的智者不要在我们面前嘲弄，除非他们自己内部先达成一致。因为他们中间有些人说有无穷的天和世界。当严密的论证推翻了他们的愚蠢体系，当几何学定律确立了根据天的本性不可能有两个天时，我们只会更加嘲笑这种精雕细琢的科学琐碎。这些博学之士看到同一原因不仅形成一个气泡，而是多个气泡，却怀疑创造智慧的能力能带来多个天！然而，如果我们把目光投向天主的全能，就会发现天的力量与伟大不同于从泉水表面冒出的水滴。因此，他们关于不可能性的论证是何等可笑！就我而言，我岂止相信第二个天，我甚至寻求第三个天，蒙福的保禄曾得以目睹。圣咏作者在说\Quote{天外之天}时，岂不暗示了天的多重性吗？天的多重性岂不比几乎全体哲学家都同意的七个行星所经过的七个轨道更为奇特？他们向我们描绘这些轨道如同桶套桶般彼此连接。他们说，这些轨道被带向与世界相反的方向，撞击以太，发出甜美和谐的声音，任何最甜美的旋律都无法比拟。如果我们请他们提供感官的证据，他们会说什么？他们说，我们从出生就习惯了这个噪音，因为一直在听，反而失去了对它的感知；就像铁匠铺里的人耳朵被不断的噪音震聋了一样。如果我驳斥这种巧妙的轻浮言论，其虚假性从第一句话就显而易见，那将显得我不知道时间的价值，并且不信任如此听众的智力。\par

但让我把外人的虚妄留给那些外人，回到教会本有的主题。如果我们相信我们之前的一些人，这里并不是创造了一个新天，而是对第一个天的新叙述。他们给出的理由是，前面的叙述简要描述了天地的创造；而这里的圣经更详细地说明了每个受造物的方式。然而，我既然圣经给这第二个天取了另一个名称并赋予其特有的功能，便坚持它与起初所造的天不同；它具有更强的本性，并对宇宙有特殊的用途。\par

4. \BibleQuote{\Emph{天主说：\Quote{在水与水之间要有穹苍，将水分开。}天主造了穹苍，把穹苍以下的水和穹苍以上的水分开了。}} \BibleRef{创 1:6} 在把握圣经的含义之前，让我们先设法应对来自其他方面的异议。有人问我们，如果穹苍是一个球形体（正如肉眼所见），其凸面如何能容纳在更高区域流动和循环的水？我们当如何回答？只有一件事：因为一个物体的内部呈现完美的凹面，并不必然意味着其外表面是球形且光滑的圆。看看浴室中的石拱顶，以及洞穴式建筑的结构；形成内部的圆顶并不阻止屋顶通常具有平坦的表面。那么，让这些不幸的人停止折磨我们和他们自己，关于我们在更高区域保留水的不可能性吧。\par

现在我们必须论及穹苍的本性，以及为何它受命位于水之间的中间位置。圣经常以\Quote{穹苍}一词表示非凡的力量。\Quote{主是我的穹苍和避难所。} \Quote{我加强了它的支柱。} \Quote{你们要在祂能力的穹苍中赞美祂。} 外教作者称密实而充满的物体为强固的物体，以区别于数学意义上的物体。数学意义上的物体仅存在于三维之中：宽、深、高。相反，强固的物体则对这些维度增加了抵抗力。圣经的习惯是将一切刚强不屈的事物称为穹苍。它甚至用这个词来表示空气的凝聚：它说，那位加强雷声者。圣经藉加强雷声意指风的强度与抵抗力，风被封闭在云层的空洞中，当它猛烈冲出时便产生雷声。因此，依我看来，穹苍是一种坚固的物质，能够保留流动而不稳定的水元素；而且，按照普遍的理解，穹苍似乎来源于水，我们不可相信它类似于冻结的水或任何其他由水渗透所产生的物质，例如水晶，据说其变形是由于过度凝结，或在矿中形成的透明石。这种透明石，若在其自然完美状态下发现，内部无裂纹，也无丝毫腐败斑点，其清晰度几乎可与空气媲美。我们不能将穹苍与这些物质中的任何一种相比。对天体持有这样的看法是幼稚而愚蠢的；尽管万物皆在万物之中：火在土中，气在水中，其他元素彼此相含；尽管我们感官所及的物事没有一种是纯粹而不混合的，不是与作为媒介的元素混合，就是与相反的元素混合；然而，我不敢断言穹苍是由这些单纯物质之一或它们的混合物形成的，因为圣经教导我不要让想象力驰骋太远。但我们要不忘指出，在这些神圣的话语\Quote{要有穹苍}之后，并没有说\Quote{穹苍就造成了}，而是\Quote{天主造了穹苍，并将水分开。} \BibleRef{创 1:7} 聋子啊，听吧！瞎子啊，看吧！——那么，谁是聋子？就是听不见圣神这令人惊愕之声的人。谁是瞎子？就是看不见独生子如此明证的人。\Quote{要有穹苍。}这是首要和本原之因的声音。\Quote{天主造了穹苍。}这是对天主主动而创造的德能的见证。\par

5. 但让我们继续解释：\Quote{\Emph{让水与水分离}}。\BibleRef{创 1:6} 从四面八方漫溢大地并悬在空中的水团是无限的，以致与其他元素不成比例。因此，如前所述，深渊覆盖了大地。我们给出这丰富水量的原因。你们中必定无人会攻击我们的意见，即使是那些最有修养、目光锐利能穿透这易朽易逝的自然的人；你们不会指责我提出不可能或空想的理论，也不会问我这流动元素以什么为基础。正因同样的理由使它们将比水更重的土地从世界两端吸引并悬浮在中心，他们无疑会同意我们：这巨大的水团静止地停留在地上，既因其向下的自然引力，也因其整体平衡。因此，庞大的水团遍布大地四周；与大地不成比例且远比它大，这是至高工匠的先见之明，祂从起初预见了将要发生的事，并在最初为世界未来的需要预备了一切。但为何需要如此超量的水？火的本质对世界是必要的，不仅对于地上出产的经营，也为了宇宙的完满；因为如果最强大、最充满活力的元素缺失，宇宙就不完美了。火与水彼此敌对并相互毁灭。火若更强，便消灭水；水若更丰沛，则消灭火。因此，必须避免这两个元素之间的公开争斗，以免因其中之一彻底消失而导致宇宙解体，至高安排者创造了如此大量的水，以致尽管火不断消耗，它仍能持续到世界毁灭的预定时刻。祂以重量和尺度筹划一切，祂按照约伯的话知道雨滴的数量，祂知道祂的工程会持续多久，以及祂应当允许火消耗多少。这就是创造时水量丰富的原因。再者，没有谁如此脱离生活，以至于需要学习火为何对世界必不可少。不仅所有维持生命的技艺——纺织、制鞋、建筑、农业——都需要火的帮助，而且树木的生长、果实的成熟、陆生和水生动物的繁殖及其滋养，都从一开始就来自热，并一直靠热的作用维持。因此，热的创造对受造物的形成和保护不可或缺，而在火持续而不可避免的消耗面前，水的丰沛也同样不可或缺。\par

6. 请你观察受造界；你将看到热力统御一切生灭之物。由于它，遍及全地的一切水，以及那些超出我们视线、散布在地底深处的水，才得以存在。由于它，才有了丰沛的泉源、溪流或井，以及江河的奔流——既有山洪也有长流不息的河流，好在许许多多不同的水库中储存湿气。据地理学家说，从东方、从冬至点流出的\Place{印度河}是地上最大的河流。从东方的中部流出\Place{巴克特鲁斯河}、\Place{科阿斯佩斯河}与\Place{阿拉克塞斯河}，\Place{塔奈斯河}从中分出，流入\Place{麦奥提斯沼泽}。再加上从\Place{高加索山}流出的\Place{法西斯河}，以及从北方地区流入\Place{攸克辛海}的无数其他河流。从西方暖国、从\Place{比利牛斯山}麓流出\Place{塔尔提苏斯河}与\Place{伊斯特河}，其中一条注入\Place{柱石之外的海}，另一条流经欧洲后注入\Place{攸克辛海}。还需要列举那些由\Place{里菲山脉}在\Place{西徐亚}腹地涌出的河流、\Place{罗讷河}以及许多其他可通航的河流么？它们灌溉了\Place{西高卢人}、\Place{克尔特人}及其附近蛮族之地，然后流入\Place{西方大海}。还有从南方更高地区流经\Place{埃塞俄比亚}的河流，有的注入我们的海，有的注入不可通行的海——\Place{厄贡河}、\Place{尼塞斯河}、\Place{克雷梅特河}，而尤其是\Place{尼罗河}，当它像海一样淹没埃及时，便不具河流的特性。如此，我们既居的大地被水环绕，通过浩瀚的海相连，并由无数常年不息的河流灌溉——这是那位安排万有的天主不可言喻的智慧，为的是防止这敌对火元素的水被彻底消灭。\par

然而，时候将到，万物都要被火吞噬；正如依撒意亚论及宇宙的天主时所说的：\BibleQuote{对深渊说：\InnerQuote{你要干涸，我必使你的河流干涸。}}\BibleRef{依 44:27}因此，你要弃绝这世间愚蠢的智慧，与我一同接受更为简单却绝不错误的真理教导。\par

7. 因此我们读到：\BibleQuote{\Emph{在水中间要有穹苍，将水与水分开。}}我曾说过圣经中\Quote{穹苍}一词的含义。它实际上并非一种坚固密实的、有重量和阻力的实体；否则这名称更适合于地。但由于在其上方的物体是轻飘、不坚固、且非我们任何感官所能把握的，是与这些纯粹而不易察觉的实体比较而言，穹苍才得了这个名称。试想一个适宜分隔湿气的地方，让纯净过滤后的湿气升入高处，而浓稠含土的湿气下降；这样，通过逐渐抽去湿气微粒，从始至终保持同样的温度。你不相信这巨量的水，但你却没有考虑到巨量的热——按体积说无疑较小，但若视为湿气的毁灭者，则异常强大。它吸引周围的水汽，正如甜瓜所示，并且一经吸引便迅速消耗，犹如灯焰吸取灯芯提供的燃料并烧尽。谁会怀疑以太是一团炽热的火？如果造物主没有为它设定不可逾越的界限，什么能阻止它点燃并烧尽近旁的一切，并吸尽存在之物中的所有湿气？是那由河流、泉源、沼泽、湖泊和海洋发出的蒸汽所形成的空中的水，遮蔽了天穹，阻止以太侵袭并烧毁宇宙。如此，我们甚至看见这个太阳在夏季瞬时晒干一片潮湿多沼泽的田野，使之完全干涸。所有的水到哪里去了？让这些无所不知的师傅告诉我们吧！难道不是每个人都明白，水已化为蒸汽并被太阳的热力消耗了吗？然而他们却说太阳甚至没有热。他们浪费了多少时间在言辞上！请看他们依靠什么证据来抵抗如此明显的事实：太阳的颜色是白的，既非淡红也非金黄。因此它不是火性的；而它的热，他们说，是由于旋转的速度。他们得到了什么？就是太阳似乎不吸收湿气？我不拒绝这个说法，虽然它是假的，因为它有助于我的论证。我说过热力的消耗需要这样巨量的水。太阳的热是本性的，还是由它的作用产生的，这没有分别，只要它对同一物质产生相同效果。如果你用两块木头摩擦生火，或者把它们靠近火焰点燃，效果完全相同。此外，我们看到那位治理万有的伟大智慧使太阳从一个区域旅行到另一个区域，以免它一直停留在同一处，其过度热力会破坏宇宙秩序。现在它在冬至前后进入南方区域，现在又回到春分点；从那里它又在夏至时前往北方区域，并通过这种不易察觉的移动，保持整个世界舒适的温度。\par

让这些博学之人看看他们是否彼此不一致。他们说，太阳消耗的水正是防止海水上涨并泛滥河流的水；太阳的温暖留下了盐分和水的苦味，而从水中吸收了纯净可饮的微粒，这是由于这一星体独特的德性——它能吸引一切轻质的，而让所有浓稠含土的东西像泥浆沉淀一样落下。由此产生了海的苦味、咸味以及枯萎和干燥的特性。而众所周知，他们坚持这些观点，却又改变立场，说湿气不能因太阳而减少。\par

8. \BibleQuote{\Emph{天主将穹苍称为天。}} \BibleRef{创 1:8} 正权属于另一者，穹苍只因与天相似而分有它。我们常发现可见区域被称为天，因我们视野内的空气稠密而连续，其名称\Quote{天}来源于意指观看的词语。圣经论及此说：\BibleQuote{空中的飞鸟}，\BibleQuote{飞鸟……在空中（穹苍）飞翔}；\BibleRef{创 1:20} 此外又说：\BibleQuote{他们升到天上。} \Person{梅瑟}祝福\Person{若瑟}支派时，为其求天的果实与露珠、夏天的太阳与月亮的交汇，以及来自山巅与永恒山丘的祝福，总之，一切滋养大地之物。在诅咒\Place{以色列}时又说：\BibleQuote{你头上的天要变为铜}。\BibleRef{申 28:23} 这是什么意思？它威胁要带来完全的干旱，缺乏那使大地果实生发的空中之水。\par

既然圣经说露水或雨水从天而降，我们便知道那是来自那些被指派占据高处的众水。当地上的蒸汽在空气的高处聚集，在风的压力下凝结时，这种空中的湿气便散布为水汽和轻云；然后再次混合，形成水滴，被自身重量拖曳而下，这就是雨的起源。当水被风的暴力击打，变成泡沫，并因极度寒冷而结冰时，它冲破云层，落下为雪。你以此可以解释空气悬浮在我们头顶的所有潮湿物质。\par

不要让任何人将哲学家对天上的好奇讨论与圣神启示的单纯与质朴相比较；正如贞洁妇女的美貌胜过妓女，我们的论证也优于对手。他们只寻求以强辩说服人。在我们这里，真理赤裸而无伪地呈现自身。但为何要自我折磨去驳斥哲学家的错误呢？只要拿出他们相互矛盾的书，作为安静的旁观者，观看这场战争就够了。因为那些思想家，在人数上并不更少，名声也不更逊，言语也不更克制，他们彼此争斗，有人说宇宙被火吞噬，又从焚毁世界的灰烬中残留的种子，万物再次复活。因此，在世界中有毁灭与重生，直至无穷。所有这一切，同样远离真理，各自找到误入歧途的旁门。\par

9. 但就水的分离而言，我不得不与教会中某些作者的意见争辩，他们在崇高概念的荫庇下，陷入了隐喻，只在水上看到了一个象征，用以表示精神的和无形的力量。在高层区域，穹苍之上，居住着更好的；在低层区域，土地和物质是恶者的居所。因此，他们说，天主由天上之水受到赞美，即是说，由善的力量，其灵魂的纯洁使他们值得歌颂天主。而天下之水则代表恶灵，他们从自然的高度堕入邪恶的深渊。狂暴、叛乱、被激情的汹涌波涛所激动，他们被称为海，因其运动的不稳定和反复无常。让我们拒绝这些理论，视之为梦呓和老妇之谈。让我们明白，水就是指水；对于穹苍对水的划分，我们接受已给予我们的理由。然而，虽然天上的水被邀请荣耀宇宙之主，我们不可将其看作有智力的存在；天并非活着，因为它们\BibleQuote{宣讲天主的荣耀}，穹苍也不是有感觉的存在，因为它\BibleQuote{显示祂的化工}。如果他们对你说，天指默观的力量，穹苍指生产善的主动力量，我们欣赏这理论之巧妙，却不能承认其真实。因为如此，露、霜、寒、热——在\BibleBook{达尼尔}{书}中被命令赞美万有创造者——将变成有智力的、不可见的本性。但这只是明智者所接受的比喻，以完成创造者的荣耀。此外，天上的水，这些因其自身拥有的德能而受特恩的水，并非唯一赞美天主的众水。\BibleQuote{你们龙和一切深渊，请从地上赞美上主。} 因此，\BibleBook{圣咏}{集}的作者没有拒绝那些我们譬喻发明家归类于邪恶领域中的深渊；他接纳他们进入宇宙的合唱团，深渊用它们的语言向创造者的荣耀唱出和谐的赞歌。\par

10. \Quote{\Emph{天主看了认为好}。} 天主并非以眼睛的悦目来判断祂工作的美，祂对美的概念也与我们不同。祂所认为的美，是那在完美中呈现艺术的一切适宜性，并趋向其目的之有用性。因此，那位在祂作品中自始便怀有明确设计者，赞许每一作品，视其按照祂创造的目的完成了自身的功用。一只手、一只眼或雕像的任何一部分，若与整体分离，无人会觉得美。但若各部分各归其位，那几乎未被察觉的比例之美，便会打动甚至最粗俗的人。但艺术家在组合作品各部分之前，就辨别并认出每一部分的美，思考他所追求的目标。圣经正是这样向我们描绘了至高艺术家，赞扬祂的每一件作品；不久，当祂的作品完成时，祂将给予整体应得的赞美。让我在此结束第二天讲道，好让勤勉的听众思考他们刚才听到的内容。愿他们的记忆为了灵魂的益处而保留；愿他们通过仔细的默想，消化我所讲的并从中获益。至于那些靠劳作生活的人，让我允许他们整天忙于自己的事务，以便他们能以无忧无虑的灵魂在晚上来赴我讲道的筵席。愿天主——在创造了如此伟大的事物之后，把这样软弱的话语放在我口中——赐予你们认识祂真理的智慧，使你们能从可见之物上升到那不可见的存在者，并使受造物的伟大与美丽使你们对创造者有正当的认识。因为自从创造世界以来，祂那看不见的事理，都可凭祂所造的万物，辨认洞察出来，以致祂的永能和神性是明明可知的。同样，地、空气、天、水、昼、夜，一切可见之物，都提醒我们谁是我们的恩主。因此，若我们借着不间断的纪念，使天主常居我们心中，我们就不会给犯罪提供机会，也不会给敌人留地步；愿一切光荣和崇拜归于祂，从今直到永远，世世无尽。阿们。\par

\SourceRecord{Translated by Blomfield Jackson.；Nicene and Post-Nicene Fathers, Second Series；Vol. 8.；Edited by Philip Schaff and Henry Wace.；Buffalo, NY: Christian Literature Publishing Co.,；1895.；<http://www.newadvent.org/fathers/32013.htm>.}

% structure: p0001 source=h1 target=section reason=page-title

\section{六日创造（第四篇讲道）}

\Emph{论水的汇聚。}\par

1. 有些城邑的居民，从早到晚，以观看无数变戏法者的把戏为乐。他们从不厌倦听那些放荡的歌曲，这些歌曲使他们的灵魂生出许多不洁；他们常被称为有福，因为他们忽略了生意的烦恼和有益于生活的行业，将他们被赐予此世的时间浪费在闲散和享乐中。他们不知道，充满不洁景象的剧场，对坐在那里的人而言，是一所共同的邪恶学校；那些悦耳而妖冶的歌曲渗入人的灵魂，所有听到的人，由于渴望模仿琴师和笛手的音调，便充满了污秽。另有一些人，对马匹痴迷，以为自己在梦中骑马；他们套上战车，更换车夫，甚至在睡眠中也不得脱离白日的愚妄。而我们，主，伟大奇工的创造者，召叫我们来默观祂自己的工程，难道我们要厌倦观看它们，或迟于聆听圣神的话语吗？我们岂不该围绕着神圣创造的广阔而多样的工场，在心中回到往昔的时光，观看创造的整个秩序？如先知的话所言，天像穹顶一样悬起；地，这巨大的质量，依靠自身而存在；环绕它的空气，柔软而流动的本质，是一切呼吸之物真实而不断的滋养，如此稀薄，以至于它稍一移动身体便退让开启，不抵抗我们的动作，同时，在瞬间，它又流回原处，在我们身后归位；最后是水，它供给人饮用，或为我们的其他需要而设计，以及它奇妙地汇聚到指定给它的确定地方：这就是我刚读过的话语将要展示给你们的景象。\par

2. \BibleQuote{ \Emph{天主说：天下诸水应聚在一处，使旱地出现。事就这样成了。} } 于是天下诸水聚在一处；\BibleQuote{天主称旱地为地，称水的汇聚为海。} \BibleRef{创 1:9} 在我先前的讲道中，你们问我为何地是看不见的，为何一切物体天然具有颜色，为何一切颜色都处于视觉之下，你们给我添了多少麻烦啊！也许，当我告诉你们说，地并非天然不可见，只是因为完全覆盖它的水层而对我们不可见时，我的理由在你们看来并不充分。那么，请听圣经如何解释自己：\BibleQuote{让水聚在一处，让旱地出现。} 幔子被揭开，让迄今不可见的地得以看见。也许你们会向我提出新的问题。首先，水往下流难道不是自然规律吗？为何圣经将此事归因于创造者的命令呢？只要水被铺展在平坦的地面上，它就不流动；它是静止的。但一遇到斜坡，前面的部分立即落下，后面的部分随之而来，第三部分又接替它。如此不停地流动，它们一个压着一个，其流速与所载的水量及所经的坡度成正比。如果水的本性如此，那么命令它聚在一处就是多余的；由于它天然的不稳定性，它必然要落入地的最凹陷处，直到其表面平整为止。我们看到，没有什么比水面更平坦的了。此外，他们还说，当我们看到几个海彼此隔着最远的距离时，这些海如何接到命令聚在一处呢？对于第一个问题，我回答说：自从天主的命令以来，你完全知道水的运动；你知道它是不稳定、不安定的，并自然地沿着斜坡和凹陷处流下。但在这一命令使其开始运行之前，它的本性是什么？你自己不知道，也没有从任何目击者那里听到。实际上，你要知道，是天主的一句话造成了本性，这一命令对受造物而言是其未来进程的指引。白天和黑夜只有一次创造，而从那一刻起，它们就不断地相互接替，平分时间。\par

3. \BibleQuote{让水聚在一处。} 它被命令成为水的自然属性，即流动，并且服从这一命令，水在它的流程中永不疲倦。我这样说，只着眼于水的流动属性。有些水天然流动，如泉水和河流；另一些则汇聚而静止。但我现在说的是流动的水。\BibleQuote{让诸水聚在一处。} 当你站在一股涌出大量泉水的泉源旁时，你难道从未想过：是谁使这水从地心涌出？是谁迫使它上升？它的储藏库在哪里？它流向何处？为什么在这里永不枯竭，在那里永不泛滥？这一切都来自那最初的命令；它对水而言是它流动的信号。\par

在水的一切叙述中，要记住这第一个命令：\Quote{让水聚在一处。}为了到达分配的位置，它们不得不流动，而一旦到达，就留在原位，不再前进。因此，用\BibleBook{训道}{篇}的话说：\BibleQuote{所有河流都奔向大海，大海却不曾满溢。}\BibleRef{训 1:6}水因天主的命令而流动，大海也根据这第一条律法被限制在界限内：\Quote{让水聚在一处。}为防止水越过河床，一次又一次地侵袭并逐一淹没所有陆地，最终淹没整个大地，它接受了聚在一处的命令。因此，我们常看到狂暴的大海掀起巨浪冲向天空，但一旦触及海岸，就在泡沫中平息其猛烈，而后退却。\BibleQuote{你们不惧怕我吗？——主说——我曾以沙粒作为海的界限。}\BibleRef{耶 5:22}一粒沙，最脆弱的东西，却遏制了海洋的暴力。因为，若不受造物主命令的约束，什么能阻止红海侵入地势较低的整个埃及，并与冲刷其海岸的另一海相连呢？若我说埃及低于红海，这是因为每次有人试图将埃及海与红海所属的印度洋连接起来时，经验都使我们确信这一点。因此，我们放弃了这一事业，同样，设想此事的埃及人\Person{塞索斯特里斯}以及后来希望实现的玛待人\Person{达理阿}也放弃了。\par

我报告这一事实，为使你们理解这命令的全部力量：\Quote{让水聚在一处}；也就是说，不再有别的聚集，并且一旦聚集，就不许它们四散。\par

4. 说水聚在一处，表明它们先前分散在许多地方。高山被深谷切割，在山谷中积蓄了水，这时水从四面八方流向那唯一的聚集处。何等广阔的平原，其范围如同广阔的大海，何等幽谷，何等以不同方式挖空的洞穴，当时都充满了水，却因天主的命令而变空！但我们不可因此说，既然水覆盖了地面，那么后来容纳海洋的所有盆地最初就是满的。如果盆地已经满了，水的聚集从何而来？我们回答说，这些盆地只在需要水汇聚成单一水体时才被预备。那时，\Place{加底拉}之外的海以及令航海者畏惧的广阔海洋，即环绕\Place{不列颠}岛和\Place{西班牙}西部的海洋，还不存在。但是，天主突然创造了这广阔的空间，水体就流了进去。\par

现在，如果我们对创造世界的解释似乎与经验相反（因为很明显，所有水并没有都流到一处），仍可作出许多回答，一旦说出就显而易见。或许甚至回答这些异议都是可笑的。他们难道要提出池塘和雨水积聚来反对，并认为这足以推翻我们的推理？显然，主要且最完全的水的汇合才被称为\Quote{聚在一处}。因为井也是水聚集之处，由人手制造，以接受散布在地洞中的湿气。\Quote{聚集}这个名称并非指任何随意的水的积聚，而是指最大且最重要的，其中元素被显示为汇集在一起。正如火，尽管被分成足以供我们这里使用的微小粒子，却在天际中成团散布；正如空气，尽管同样微小地分划，却占据了大地周围的区域；同样，水，尽管少量地四处散布，只形成一个聚集，即把整个元素与其他部分分开。无疑，湖泊，无论是北方的那些，还是在\Place{希腊}、\Place{马其顿}、\Place{比提尼亚}和\Place{巴勒斯坦}发现的，都是水的聚集；但这里指的是最大的，其范围与大地相等。前者含有大量的水；无人否认。然而，没有人能合理地称它们为海，即使它们像大海一样充满盐和沙。例如，\Place{犹太}的\Place{阿司法提提斯湖}，以及\Place{阿拉伯沙漠}中位于\Place{埃及}和\Place{巴勒斯坦}之间的\Place{色尔波尼湖}。这些是湖泊，只有一个海，正如环游大地的人所确认的。尽管一些权威认为\Place{希尔卡尼亚海}和\Place{里海}是封闭在自己的界限内，但如果我们相信地理学家，它们彼此相通，并共同注入\Place{大海}。因此，根据他们的说法，\Place{红海}和\Place{加底拉}之外的海只形成一个。那么，为什么天主称不同的水体为海呢？原因是：水流到一处，而它们的不同积聚，即大地在其褶皱中所拥抱的海湾，从主那里得到了海的名字：\Place{北海}、\Place{南海}、\Place{东海}、\Place{西海}。海甚至有自己的名字：\Place{攸克辛海}、\Place{普罗庞提斯海}、\Place{赫勒斯滂海峡}、\Place{爱琴海}、\Place{爱奥尼亚海}、\Place{撒丁尼亚海}、\Place{西西里海}、\Place{第勒尼安海}，以及许多其他名字，现在列举它们会太长，而且不太合适。看，为什么天主称水的聚集为海。但让我们回到我的论证所偏离的出发点。\par

5. 天主说：\BibleQuote{\Emph{让水聚集在一处，让旱地出现}}。祂没有说\Quote{让地出现}，免得它再次显示为无形、泥泞并与水混合，也不具有适当的形状和能力。同时，以免我们将旱地的形成归功于太阳，造物主在太阳受造之前就向我们显示它已干涸。让我们跟随圣经所给予的思想。不仅覆盖大地的水从它流走，而且所有渗入其深处的水也因着至高主宰不可抗拒的命令而退去。事就这样成了。这足以表明造物主的声音产生了效果；然而，在几种版本中，还加上了\Quote{水聚到了一处，旱地显露出来}——其他译者没有给出这些话，它们似乎不符合希伯来用语。事实上，在断言\Quote{事就这样成了}之后，重复完全相同的话是多余的。在准确的抄本中，这些话被标上了删节号，这是删除的记号。\par

\BibleQuote{\Emph{天主称旱地为地，称水的汇聚为海}}。\BibleRef{创 1:10}为何圣经先说\BibleQuote{水聚集在一处，旱地出现}？为何这里又加上旱地出现，而天主给它起名为地？这是因为干燥似乎是表征主体本性的属性，而\Quote{地}这个词只是它的简单名称。正如理性是人的特有官能，\Quote{人}这个词用来指称拥有这一官能的存在者，同样，干燥是地特殊而独有的性质。因此，本质干燥的元素就被称为地，正如以嘶叫为特征叫声的动物被称为马。其他元素，如同地一样，都接受了某种特殊属性，使它们与其余元素区别开来，并表明它们是什么。因此，水的特性是冷，空气是湿，火是热。但这一理论实际上只适用于世界的原始元素。那些参与构成物体并为我们感官所感知的元素，向我们展示了这些性质的结合；在全部自然中，我们的眼睛和感官找不到任何完全单一、简单和纯粹的东西。地同时是干和冷；水是冷和湿；空气是湿和温；火是温而干。正是通过它们性质的结合，不同的元素才能混合。由于共同的属性，每个元素都与邻近元素混合，而这种自然的联盟将它连接到相反的元素。例如，地既干又冷，在冷中找到了与水相连的关系，并借助水与空气相连。水被置于两者之间，似乎与两者都握手，凭借其双重性质，通过冷与地结盟，通过湿与空气结盟。空气则占据中间位置，在水与火的敌对本性之间扮演中介角色，通过湿与第一个结合，通过热与第二个结合。最后，火的本性既温又干，通过温暖与空气相连，并通过干燥与地重新结合。从这种和谐与元素的相互混合中，产生了一个圆形的和谐合唱，每个元素都因此得名。我说这些是为了解释为何天主将旱地称为地，却没有将地称为干。这是因为干燥不是地后来获得的那些性质之一，而是从一开始就构成其本质的性质。使一个物体存在的原因，自然先于它后来的性质，并具有凌驾于它们之上的优位。因此，天主合理选择了地最古老的特征来称呼它。\par

6. \BibleQuote{\Emph{天主看了认为好}}。\BibleRef{创 1:10}圣经并非仅仅想说大海呈现给天主的悦目景象。造物主并非用眼睛观看祂工程的美丽。祂是在祂不可言喻的智慧中默观它们。大海在平静安宁中一片明亮，是美景；当被一阵微风轻拂，海面显出紫色和蔚蓝的色彩，它没有猛烈拍打邻近的海岸，而是似乎用平静的爱抚亲吻它们，这也同样是美景。然而，圣经并非由此让天主找到大海的良善与魅力。这里，工程的目的才造就了良善。\par

首先，海水是地上一切湿气的来源。它通过不可见的管道渗透，正如波浪所进入的地下洞穴所证明的；它被接收在倾斜蜿蜒的渠道中；然后，被风驱动，它上升到地面并破土而出，经过长时间的渗滤变得可饮用且去除了苦味。常常由于同样的原因，它甚至从所经过的矿脉中涌出，从中获得温暖，并沸腾升起，喷出炽热，正如在岛屿和海岸上所见的；甚至在远离海洋的内陆某些地方，在河流附近，以小事比大事，几乎同样的现象也会发生。这些话意味着什么？是为了证明大地底下全是由看不见的渠道挖空的，水从海洋源头在底下到处流动。\par

7. 如此，在天主眼中，海是善的，因为它使地底深处保持湿气的地下水流。它又是善的，因为它从四面八方接受河流却不超越自己的界限。它是善的，因为它是空中水气的源头和来源。被太阳光线加热后，它化为蒸汽，被吸引到高空区域，在那里因升得太高而超出地面光线的折射而被冷却，再加上云影的冷却作用，它变为雨水，使大地肥沃。若有人不信，请看火上的锅，虽然装满水，却常因水全部煮沸化为蒸汽而变空。水手们甚至煮沸海水，用海绵收集蒸汽，以解燃眉之急。\par

最后，在天主眼中，海是善的，因为它环绕岛屿，同时构成它们的屏障与美丽；因为它将大地最遥远的部分连结起来，促进水手之间的相互交流。借此，它赐予我们普遍信息的恩惠，供给商人财富，并轻松提供生活必需品，使富人得以输出剩余，使穷人得以获得所缺的供应。\par

然而，我从何处感知到海洋在造物主眼中的美善呢？若海洋在天主面前是善而可赞的，那么像这样一座教会的集会——其中男人、妇女和儿童的声音在祈祷中向天主升起，如拍岸波涛般混响——该是多么更美好啊！这座教会也享有一片深邃的宁静，恶灵不能以异端的气息扰乱它。那么，愿你们忠于此等善导，以获得主的赞许，在我们的主耶稣基督内，愿光荣与权能归于祂，至于无穷之世。阿们。\par

\SourceRecord{Translated by Blomfield Jackson.；Nicene and Post-Nicene Fathers, Second Series；Vol. 8.；Edited by Philip Schaff and Henry Wace.；Buffalo, NY: Christian Literature Publishing Co.,；1895.；<http://www.newadvent.org/fathers/32014.htm>.}

% structure: p0001 source=h1 target=section reason=page-title

\section{\WorkTitle{Hexaemeron}（第五篇讲道）}

\Emph{大地的发芽}。\par

1. \BibleQuote{\Emph{天主说：\Quote{地要生出青草，结种子的菜蔬，并结果子的树木，各从其类，果子都包着核。}}} \BibleRef{创 1:11} 那是深奥的智慧命令大地，在卸下水的重压而安息之后，首先长出青草，然后长出树木，正如我们现在仍看到它此时所做的。因为那时所听到的声音和这道命令，对大地来说，成了自然而永久的法则；它赐予了丰产和在未来所有时代结出果实的能力；\BibleQuote{地要生出。} 植物的出产首先显示出发芽。当幼芽开始萌发，它们形成青草；这青草发育并成为植物，不知不觉地获得其不同的关节，并在种子中达到成熟。这样，所有发芽和青绿的东西都发展了。\BibleQuote{地要生出青草。} 大地要凭自己生出，无需任何外来的帮助。有些人认为太阳是大地一切生产力的来源。他们说，是太阳热量的作用将生命力从地心吸引到地表。大地装饰在太阳之前的原因如下：为使那些崇拜太阳为生命之源的人，能够放弃他们的错误。如果他们确信大地在太阳生成之前就已装饰好了，他们就会收回对太阳的无限崇拜，因为他们看到青草和植物在太阳升起之前就已生长。如果为羊群预备了食物，难道我们的种族就不值得同样关怀吗？那为马和牛提供牧草的，首先想到的是你的财富和享乐。他喂养你的牲畜，是为了供应你生活的一切需要。而五谷的生出，若不是为了你的生存，又是为了什么目的？此外，许多草类和蔬菜也作为人的食物。\par

2. \BibleQuote{\Emph{地要生出结种子的菜蔬，各从其类。}} 因此，虽然某些种类的草对动物有益，但它们的收益也是我们的收益，而种子特别为了我们的用途。这就是我所引用的话的真正含义。\BibleQuote{地要生出青草，结种子的菜蔬，各从其类。} 这样我们可以恢复这些话的顺序，在现行译本中其结构似乎有缺陷，而自然的法则将得到严格遵守。事实上，首先是发芽，然后是青绿，接着是植物的生长，在达到完全生长后，在种子中达到完美。\par

那么，他们怎么说，圣经能否描述大地所有的植物都结种子，而芦苇、匍匐冰草、薄荷、番红花、大蒜、花蔺以及无数其他种类却不结种子？对此我们回答说，许多蔬菜的生殖力在根部和下方。例如，竹子每年生长后从根部发出一个突起，代替未来的种子。许多其他蔬菜也是如此，它们都通过根在地球上繁殖。没有比每种植物都产生自己的种子或含有某种生殖力更真实的了；这就是所谓\BibleQuote{各从其类}的意思。所以，芦苇的芽不会产生橄榄树，而是从芦苇长出另一棵芦苇，从一种种子总是长出同一类的植物。这样，所有从大地最初生出时生长的东西，直到我们的时代都保持不变，归功于种类的不断繁殖。\par

\BibleQuote{地要生出。} 看，在这短短的话语、在这简短的命令下，冰冷贫瘠的大地开始产痛，急忙结出果实，它脱去悲伤阴沉的覆盖，披上更灿烂的外衣，以自身的装饰而自豪，展现出植物的无限多样性。\par

我希望受造物如此充满你的赞叹，以至于无论你在哪里，最小的植物都能使你清晰地记起造物主。如果你看到田野的青草，要想到人的本性，并记住\Person{依撒意亚}智者的比喻。\BibleQuote{凡有血肉的尽如草，一切荣华都像野地的花。} 的确，生命的迅速流逝，短暂的满足和幸福瞬间给予一个人的快乐，都奇妙地符合先知的比喻。今天他身体健壮，因奢侈而肥胖，年富力强，面色如花，强壮有力，具有不可抵抗的精力；明天他就成为可怜的对象，因年老而枯萎或因疾病而虚弱。另一个人闪耀着辉煌财富的全部光彩，周围有一大群谄媚者，一群假朋友追逐他的恩宠；一群亲属，却不是真正的亲属；一大群仆人簇拥着他，为他的食物和一切需要服务；在他出入时，这无数随从跟随着他，引起所有遇见他的人的嫉妒。财富之上还可加上国家权力、皇位授予的荣誉、省总督或军队的指挥权；一个先驱者走在他前面，高声呼喊；左右的扈从也使他的臣民充满敬畏，鞭打、没收、流放、监禁，以及所有他用来让所统治之人感到难以忍受的恐惧的手段。然后呢？一夜之间，一次发烧、胸膜炎或肺炎，就把这个人从人群中间夺走，瞬间剥去他所有的舞台装饰，而这一切荣耀都被证明是一场梦。因此，先知将人的荣耀比作最脆弱的花。\par

3. 直到这一点，植物萌芽的次序为它们最初的安排提供了见证。每一种草本、每一种植物都出自一个胚芽。如果它像匍匐冰草和番红花那样，从其根部并从这一较低突起处发出嫩芽，则它必须总是发芽并向外生长。如果它来自一颗种子，那么仍然必然先有胚芽，然后有嫩芽，然后有绿叶，最后是在迄今干燥粗壮的茎上成熟的果实。\BibleQuote{让地上生出青草}。当种子落入含有适当热度和湿度的泥土中时，它膨胀并变得多孔，抓住周围的泥土，吸引所有适合它并与它有亲缘的东西到自身。这些土粒，无论多么微小，当它们落下并渗入种子的所有孔隙时，加增了种子的体积，并使其向下生根，向上发芽，发出与根数量相当的茎。由于胚芽不断变暖，通过根抽吸上来的水分，在热力吸引的帮助下，从土壤中吸取适量的营养，并将其分配到茎、树皮、谷壳、种子本身以及武装着它的芒刺上。正是由于这些连续的增生，每一株植物达到了其自然的发育，无论是谷物还是蔬菜、药草或灌木。一株植物，一片草叶足以占据你的全部才智，去沉思那产生它的技艺。为什么麦秆有节更好？它们岂不就像紧固件，帮助麦秆轻松承受穗的重量，当穗因果实饱满而向地面弯垂时？因此，燕麦顶端没有需要承受的重量，故没有这些支撑；而大自然为小麦提供了它们。它将谷粒藏在一个壳中，使其免遭鸟类的劫掠，并用芒刺的壁垒武装它，这些芒刺像标枪一样保护它免受微小生物的袭击。\par

4. 我该说什么？我该留下什么不说？在创造的丰富宝藏中，难以选择最珍贵的东西；省略之损失过于严重。\BibleQuote{让地上生出青草}；顷刻间，伴随着有用的植物，出现了有害的植物；伴随着谷物，出现了毒参；伴随着其他营养植物，出现了菟葵、附子、曼德拉草和罂粟汁。那么，我们难道不该为如此众多的有益恩赐表示感谢，反而因那些可能危害我们生命的物而责备天主吗？我们难道不该反思，并非一切都是为了我们肚腹的需要而造的？滋养的植物，注定为我们所用，近在咫尺，且为举世所知。但在受造物中，没有一样是无缘无故存在的。公牛的血是毒药：难道这头对人类如此有用的动物，本不该被创造，或者若被创造，就该没有血吗？但你自己有足够理智，使你远离致命之物。什么！绵羊和山羊知道如何避开威胁它们生命的东西，仅凭本能就能辨别危险：而你，拥有理智和医学来提供所需，还有祖先的经验告诉你避开所有危险之物，你竟告诉我你难以使自己远离毒药！但没有一样东西是无缘无故被造的，没有一样是无用的。一样东西用作某动物的食物；医学在另一样中找到了缓解我们某些疾病的方法。例如，椋鸟吃毒参，其体质使它对毒药的作用不敏感。由于它心脏孔隙的纤细，毒汁一旦被吞下，在它的寒凉攻击要害部位之前就被消化了。鹌鹑，凭借其特殊的体质，逃避了危险的效果，以菟葵为食。有时甚至毒药对人类也有用；医生用曼德拉草使我们入睡；用鸦片平息剧痛。毒参至今已被用来平息顽固疾病的暴怒；菟葵也常常消除了长期的疾病。这些植物，非但不应使你责备天主，反而给了你一个新的感恩的理由。\par

5. \Quote{\Emph{让大地生出青草}}。这些话语包含了何等自发的预备——那存在于根、植物本身和果实中的，以及我们的劳作和耕种所增添的！天主并没有命令大地立刻结出种子和果实，而是让她生出胚芽，变绿，直到种子成熟；因此，这第一个命令教导了本性在岁月进程中当行之事。但是，他们问道：大地真的按种类结出种子吗？因为常常在播种小麦后，我们却收获黑色的谷粒？这不是种类的改变，而是谷粒的变异、一种病害。它并没有不再是小麦；是因为被灼烧而变黑，正如从其名称可以得知。若严重的霜冻烧灼了它，它会有另一种颜色和不同的味道。他们甚至声称，若它能找到合适的土壤和温和的温度，它可能恢复其原形。因此，你在本性中找不到任何违反天主命令的事。至于毒麦和所有那些混杂在收成中的劣质谷粒，圣经中的莠子，远非谷物的变种，它们有自己的起源和种类；是那些篡改主之教导者的形象，他们在圣言上没有受到正确的教导，反而被邪恶者的教义所败坏，混入教会的健全身体，在更纯洁的灵魂中秘密传播他们的有害谬误。主将信他之人的完美比作种子的生长：\BibleQuote{如同一个人把种子撒在地里，他黑夜白天，或睡或起，种子发芽生长，至于怎样，他却不知道，因为大地自然结出果实：先发苗，后吐穗，最后穗上满结麦粒。}\BibleRef{谷 4:26-28} \Quote{让大地生出青草。}一瞬间，大地开始发芽，服从造物主的律法，完成了生长的每个阶段，使胚芽臻于完美。草地覆盖着深密的青草，肥沃的平原因收获而摇动，谷物的翻滚如同海浪。每棵植物、每株草、最小的灌木、最细微的蔬菜，都从大地丰盛地升起。在这最初的植被中没有失败：没有农夫的不熟练，没有恶劣的天气，没有任何事物能伤害它；那时，定罪之判词还未束缚大地的肥沃。这一切都是在那谴责我们靠额上汗水吃面包的罪之前。\par

6. \Quote{\Emph{让大地}}，造物主接着说，\InnerQuote{\Emph{结出按种类结果子的果树，其种子在它里面}}。\BibleRef{创 1:11}\par

在这命令下，每个树林都密植起来；所有的树木，如枞树、雪松、柏树、松树，都长到最高；灌木立刻披上浓密的叶子。称为冠冕植物的植物，如玫瑰、桃金娘、月桂，原本不存在；但在一瞬间它们出现了，各自带着独特的特征。最显著的区别将它们与其他植物分开，每一种都有自己独特的品性。但那时玫瑰没有刺；此后荆棘被添加到它的美丽中，让我们感到痛苦与快乐紧密相连，并提醒我们自己的罪，这罪使大地长出荆棘和蒺藜。然而，他们说，大地接受了命令要产生\Quote{种子在其中，}的果树，而我们看到许多树木既没有果子也没有种子。我们将如何回答？首先，只提到了最主要的树木；其次，仔细的考察会让我们看到每棵树都有种子，或某种代替种子的东西。黑杨、柳树、榆树、白杨，所有这类树木，都不产出明显的果实；然而，细心的观察者会在每棵树上找到种子。这种位于叶片基部的小粒——那些忙于造词的人称之为\OriginalTerm{小粒}{\GreekText{μίσχος}}——具有种子的性质。有些树通过枝条繁殖，从枝条中生出根来。也许我们甚至应该把从树根发出的幼芽视为种子：因为栽培者会拔掉它们以繁殖物种。但我们已经说过，这主要是关于那些对我们的生命贡献最大的树木；它们向人提供各种果实，并给人丰富的营养。比如葡萄树，它产酒以悦人心；比如橄榄树，它的果子以油使人面容光亮。在同一植物中，有多少事物结合！在葡萄树中，有根、嫩而柔韧的枝条（它们在地面蔓延甚远）、芽、卷须、酸葡萄和熟葡萄的串。当聪明的眼睛观察葡萄树时，它提醒你你的本性。无疑你记得那个比喻，其中主称自己为葡萄树，称祂的父为园丁，而每一个因信德被接枝到教会中的我们都是枝条。祂邀请我们结出丰硕的果实，以免我们的不结果子定我们受火刑。\BibleRef{若 15:1}祂不断将我们的灵魂比作葡萄树。祂说：\Quote{我的爱人，}\Quote{有一座葡萄园，在肥美的山冈上，}\BibleRef{依 5:1}别处又说：\Quote{我栽种了一个葡萄园，周围围上篱笆。}\BibleRef{玛 21:33}显然，祂称人的灵魂为祂的葡萄园，这些灵魂被祂诫命的权威和天使的护卫所环绕。\Quote{上主的天使在敬畏祂的人周围扎营。}再者：祂为我们栽种了支柱，可以说，在祂的教会中设立了宗徒、先知、教师；并通过古代圣者的榜样提升我们的思想，不让他们拖曳在地上被踩踏。祂希望爱的拥抱如同葡萄树的卷须，将我们与邻人相连，使我们依靠他们，这样，在不断地向往天堂时，我们就能模仿那些爬到最高树顶的葡萄树。祂也要求我们允许自己被挖掘；这正是灵魂所做的，当它摆脱世上的挂虑——这些挂虑压在我们心上。那么，谁从肉体的情感和对财富的爱中获得自由，不仅不被它们眩惑，反而藐视和轻看这可怜虚荣的荣耀，他就如同被挖掘，最终得以呼吸，摆脱了尘世思虑的无用负担。按照比喻的精神，我们也不应长出太多的木，即过奢华的生活，博取世人的喝彩；我们应结出果实，将我们工作的证明留给园丁。要\Quote{像天主殿中的青橄榄树，}永不缺乏希望，而是因信德披上救恩的花朵。这样，你将与此植物的常青相似，并在多结果实上与之争胜，若你每天大量施舍。\par

7. 但让我们回到对造化巧妙安排的考察。那时出现了多少树木，有些给我们果实，有些给我们盖房，有些给我们造船，有些给我们生火！它们各部分的结构是多么多样！然而，要找出每棵树独有的特性，并把握它们与其他物种的区别，又是多么困难！有些树扎根很深，有些则不然；有些笔直向上只有一根主干，有些似乎喜爱大地，从根部分出许多枝条。那些长枝伸向高空者，也有深根在宽阔的范围内伸展，那是自然为支撑树的重量所置的真正基础。树皮的种类多么丰富！有些树皮光滑，有些粗糙；有些只有一层，有些有多层。多么奇妙！你可以在植物的幼年和老年中找到与人类相似之处。年轻时，它们树皮膨胀；衰老时，则粗糙起皱。若砍去一棵，它发出新芽；另一棵则从此不结果，如同受了致命伤。此外，据观察，松树被砍伐甚至被火焚烧后，会变成一片橡树林。我们还知道农夫的技巧能补救某些树木的天然缺陷。例如，酸石榴和苦杏仁，若在树干近根处钻一个孔，将一块松油塞插入髓心，则失去酸汁而成为美味果实。那么，罪人也不应绝望，他想，如果农业能改变植物的汁液，灵魂努力修德，当然能战胜一切软弱。\par

果树的果实种类繁多，不可胜言；不仅不同种的果树果实多样，即使同种之内亦复如是，如果园丁所说属实——树的性别会影响果实的特性。他们在棕榈树中区分雄雌；有时我们看到所谓雌树低垂枝条，仿佛情欲炽烈，邀雄树拥抱。于是照料这些植物的人将雄棕榈树的花粉——他们称之为\Emph{psen}——撒在这些棕榈树上：树似乎分享交合的愉悦；随即枝条抬起，树叶恢复平常形态。无花果树也是如此。有的在栽培的无花果树旁种植野无花果树，另有人为补救园中结果无花果树之弱，将未熟的无花果系于枝上，以此保留已开始掉落而损失的果实。自然在此给我们何教训？即我们常常需要从信仰外人那里借取某种力量来彰显善工。若你在教会之外，在外教生活中，或在有害的异端中，见到德性与忠于道德律的榜样，就当加倍努力，效仿结果的无花果树：它借野无花果树之旁获得力量，防止果实脱落，并更细心地滋养果实。\par

8. 植物繁殖方式繁多，我们仅能略述其要。至于果实本身，谁能细数其品种、形状、颜色、特殊滋味及各各用途？为何有些果实暴露于阳光之下而成熟，而另一些却藏在壳内长大？果实柔嫩的树木，如无花果树，有浓密的叶荫；反之，果实较坚硬的树木，如胡桃，仅有稀疏的叶荫。前者娇嫩，需更多照料；后者若外壳更厚，叶荫反为有害。何以葡萄叶呈锯齿状，若非为使果穗既能抵抗空气之害，又可通过缝隙接受全部阳光？凡事皆有原因，无出于偶然。一切皆显示不可言喻之智。\par

什么言论能触及一切？人心能精确回顾，注意每种特性，展示所有差异，确切揭示如此众多奥秘的原因吗？同一水，经由根部吸取，以不同方式滋养根本身、树干之皮、木质及木髓。它变成叶，分布于枝与细枝，使果实膨大——它给予植物其树脂与汁液。谁能向我们解释这一切的差异？乳香之树脂与香脂之汁液不同，埃及与利比亚茴香所蒸馏之物亦不同。琥珀据说是植物汁液结晶而成。其明证：可见稻草碎屑与小昆虫在汁液尚液态时被捕捉而囚禁其中。一言以蔽之，无长期经验者无法找到词语表达其功效。再者，这水如何在葡萄中成为酒，在橄榄中成为油？然而令人惊奇的，并非它在一种果实中变甜，在另一种中变油润，而是在甜的果实中有不可言喻的滋味多样性。葡萄有一种甜，苹果有另一种，无花果又有另一种，海枣亦然。我愿让你们继续此研究，以此为乐。为何这水有时有甜味，因停留于某些植物而变柔和；有时则因经过其他植物而成酸，刺激味觉？又为何它达至极端苦涩，在苦艾与药牵牛中使口粗糙？在橡实与山茱萸中它有何尖锐粗糙之味？在笃耨香树与胡桃树中它变为柔软油质之物？\par

9. 但何需继续？在同一无花果树中，我们就有最相反的味道：汁液苦涩而果实甘甜。在葡萄中，岂非葡萄甜而枝条涩？色彩何其多样！看，在草地上，同一水在此花中变红，在彼花中变紫，在此花中变蓝，在彼花中变白。这色彩之多样，岂可与香气之多样相比？但我察觉，一种贪得无厌的好奇心将我的言论引至其界限之外。若我不住脚，将其召回创造之律，则白昼将尽，而我仍使你们在微小事物中看见大智慧。\par

\BibleQuote{于是地产生了果树，结有果实。} 于是山顶立刻披上绿叶：乐园精巧布局，无数植物装饰河岸。有些供人餐桌之饰，有些以其果实和树叶滋养动物，有些以其汁液、浆液、木屑、树皮或果实提供医药帮助。总之，世世代代的经验，利用一切偶然，未能发现任何有用之物，是造物主的先见之明未曾先见并呼唤存在的。因此，当你看见园中或林中的树木，喜水或喜陆的，开花的或不开花的，我愿你甚至在微小事物中认出伟大，不断增添对造物主的惊叹，并加倍爱祂。问你自己，为何祂使一些树常青，另一些落叶；为何在常青树中，有些落叶，有些则始终存留。例如橄榄与松树落叶，虽不知不觉更新，似永不失去青翠。反之，棕榈树从出生到死亡，始终以同样枝叶为饰。再想柽柳的双重生命：它是水生植物，却又覆盖旷野。因此耶肋米亚将它比作最坏的性格——双面性格。\par

10. \Quote{\Emph{愿大地生出}}。这简短的命令片刻之间便成了一种广阔的本性，一套精密的体系。它比思想更快地产生了植物那数不清的特质。正是这道命令，至今仍加诸大地，在每一年的循环中彰显它产生草木、种籽和树木的全部力量。如同陀螺，在第一次推动之后，继续旋转，一旦固定中心便绕自身转动；同样，本性接受了这第一道命令的推动，便不间断地追随岁月的进程，直到万物的终结。愿我们都急忙奔向这终结，满载果实与善工；如此，栽种在主的殿中，我们将在我们天主的庭院里繁荣，在我们的主耶稣基督内，愿光荣与权能归于祂，至于无穷之世。阿们。\par

\SourceRecord{Translated by Blomfield Jackson.；Nicene and Post-Nicene Fathers, Second Series；Vol. 8.；Edited by Philip Schaff and Henry Wace.；Buffalo, NY: Christian Literature Publishing Co.,；1895.；<http://www.newadvent.org/fathers/32015.htm>.}

% structure: p0001 source=h1 target=section reason=page-title

\section{六日创世（第六篇讲道）}

\Emph{发光体的创造。}\par

1. 在竞技场的表演中，观众必须投入运动员的努力。表演的法则指明这一点，因为它们规定所有人在竞技场中必须脱去头巾。在我看来，其目的是使每个人不仅成为运动员的观众，而且在某种程度上成为一个真正的运动员。因此，为了研究创造这一伟大而奇妙的表演，为了理解至高而不可言喻的智慧，你必须为默观我展现在你眼前的奇观而带来个人的光明，并在这场战斗中按照你的能力帮助我——你并非作为评判者，而是作为战友，以免真理逃脱你，也免得我的错误对你们共同造成损害。为何说这些话？这是因为我们打算将世界作为一个整体来研究，并思考宇宙，不是凭着世俗智慧的光，而是凭着天主启示祂仆人时所赐的光——此时祂亲自对他说话，且没有谜语。这是因为一切热爱伟大壮观表演的人都必须带来一颗充分准备好的心灵来研究它们。如果有时，在明亮的夜晚，你用警觉的眼睛凝视着星宿难以言表的美，因而想到万物的造物主；如果你曾自问是谁用这些花卉点缀了天空，为何可见之物比美丽更有益；如果有时在白天，你曾研究过光的奇迹，如果你曾藉可见之物提升自己达到不可见的存在——那么你就是一个准备充分的听者，可以在这庄严而蒙福的竞技场中就座。如同一个不认识城市的人被拉着引导穿越全城，我也要如此带领你们，像陌生人一样，游览这座宇宙大城的奇妙奥秘。我们的第一故乡曾在这座大城中，从那杀人恶魔诱惑人陷入奴役的地方，他驱逐了我们。在那里，你将看到人的最初起源，以及死亡立刻攫取他——死亡是由罪带来的，罪是恶神的长子。你将知道你是由泥土所造，却是天主双手的工程；你比牲畜软弱得多，却被命定统治无灵无理性的万物；就自然禀赋而言你处于劣势，但感谢理性的特权，你能提升自己直至天堂。如果我们被这些真理深入心底，我们就会认识自己，认识天主，敬拜我们的造物主，事奉我们的主，光荣我们的父，爱戴我们的供养者，赞美我们的恩主，我们将永不停止地尊崇今生和来世生命的主宰；祂借着今世倾注给我们的财富，使我们相信祂的应许，并用现世的美善来坚固我们对未来的期盼。确实，如果时间的美善如此，永恒的美善又将如何？如果可见之物如此美丽，我们该怎样思考不可见之物？如果天空的高远超乎人类理智的尺度，什么心智能描摹永恒者的本性？如果那受造而成、有朽坏的太阳如此美丽、如此伟大、运行如此迅疾、轨道如此恒常；如果它的宏大与宇宙如此和谐、如此恰如其分；如果因本性的美丽，它如一只明亮的眼睛闪耀在造化的中心；如果最后，人永远不会厌倦默观它——那么义德的太阳的美丽又将如何？如果盲人因看不到物质的太阳而痛苦，那么罪人不能享受真光，是何等大的损失啊！\par

2. \BibleQuote{\Emph{天主说：让天空中要有光体，照耀大地，分别昼夜。}} 天地是最先创造的；随后光被造；白昼与黑夜被分开；然后出现了穹苍和旱地。水已聚集到它被指定的蓄水池中，地显出它的产物，它使多种菜蔬发芽，并被各种植物点缀。然而，太阳和月亮尚未存在，为的是使那些不认识天主的人不致将太阳视为光的源头和父亲，或视为地上万物的创造者。因此，到了第四天，天主说：\BibleQuote{让天空中有光体。}\par

一旦你知道了是谁说话，就要立即想到听者。天主说：\BibleQuote{让天空中有光体……天主造了两个大光体。} 是谁说话？是谁造物？你岂不见一个双重位格吗？无处不在的神圣语言中，历史都播撒了神学的信理。\par

创造光体的动机随之而来：是为了照耀大地。光已经创造了，那么为什么说太阳被造是为了发光呢？首先，请不要因这表达的奇特而发笑。我们不追求你们在字词上的精细，也不甚在意使它们和谐悦耳。我们的作者不以润饰句段为乐，我们在任何地方都宁愿选择词语的清晰而非响亮。请看，神圣的作者是否是通过 \Quote{照亮} 这一表达足够表明他的思想？他将 \Quote{发光} 用作 \Quote{照耀} 的意思。这里与以前关于光所说的话并无矛盾。那时，光的实际本性被创造出来；现在，太阳的形体被构造为原始光的载体。一盏灯不是火。火有照明的特性，我们发明了灯用以在黑暗中照明。同样，光体被构造为那纯洁、清澈、非物质之光的载体。宗徒向我们谈论某些在世界中照耀的光体——它们不与世界的真光相混淆；拥有这真光使圣人们成为他们教导之人的光体，将他们从无知的黑暗中带出。因此，万物的造物主在光荣的光之外又造了太阳，并安置它在天空中照耀。\par

3. 任何人不要以为光的明亮是一回事，而作为其物质载体的物体是另一回事是不可能的事。首先，在所有复合物中，我们区分能接受性质的实体和它所接受的性质。白白的本性是一回事，被漂白的物体的本性是另一回事；这样，我们刚刚看到被造物主的能力重新联合的本性是不同的。不要告诉我分离它们是不可能的。我甚至不假装能够将光与太阳的物体分离；但我坚持认为我们在思想中分离的东西，可以被本性的造物主在现实中分离。此外，你不能将火的明亮与它所拥有的燃烧之德分离；但是天主，愿意通过一个奇妙的景象吸引祂的仆人，在燃烧的荆棘中设置了一种火，它显示了火焰的所有光辉，而其吞噬性的属性却蛰伏着。这就是圣咏作者所肯定的，他说：\Quote{主的聲音把火焰分開。} 因此，在今生之后等待我们的报应中，一个神秘的声音似乎告诉我们，火的双重本性将被分开；义人将享受它的光，而它的热的折磨将成为恶人的酷刑。\par

在月亮的圆缺变化中，我们再次找到我们所说的证明。当它停止并变小时，它不会消耗其整个身体，而是根据它沉淀或吸收其周围的光的多少，向我们呈现其减少或增加的图像。如果我们想要一个明显的证明，证明月亮在休息时不会消耗其身体，我们只需睁开眼睛。如果你在无云晴朗的天空中观察它，当它呈现完整的月牙形状时，你会发现黑暗且未被照亮的部分描绘出一个与满月形成的圆相等的圆。因此，眼睛可以看见整个圆，如果它将这个阴暗的曲线加到被照亮的部分上。不要告诉我月亮的光是借来的，随着它接近或远离太阳而减少或增加。这不是我们目前研究的目标；我们只想证明它的身体与使其发光的亮光不同。我希望你对太阳有同样的看法；不过，唯一不同的是，太阳在一次性接受光并将其与自己的物质混合后，不再放下它；而月亮则轮流脱去并重新穿上光，通过它自身所发生的事证明了我们对太阳所说的话。\par

太阳和月亮就这样接受了将白昼与黑夜分开的命令。天主早已将光与黑暗分开；然后祂将它们的本性对立，使它们不能混合，黑暗与光之间永远不会有任何共同之处。你看到白天的影子是什么；那正是黑夜的黑暗本性。如果，在光出现时，影子总是落在相反的一边；如果在早晨它伸向落日的一边；如果在晚上它倾向日出的一边，而在正午转向北方；那么黑夜就退到与太阳光线相对的区域，因为它本质上只是地球的影子。因为，就像在白日，影子是由遮挡光的物体产生的，黑夜自然地在环绕地球的空气被阴影覆盖时来临。这正是圣经所说的：\Quote{天主将光与黑暗分开。} 因此，黑暗在光来临时逃离，两者在最初创造时由于自然的相互排斥而被分开。现在，天主命令太阳测量白昼，命令月亮，每当它圆盘时，统治黑夜。因为此时这两个发光体几乎完全相对；当太阳升起时，满月从地平线消失，在太阳落山时从东方重新出现。如果在其他相位，月亮的光并不完全与黑夜对应，这对我们的主题无关紧要。尽管如此，当它达到完美时，它使星星变得苍白，并以它光的光辉照亮大地，它统治黑夜，并与太阳合作将黑夜的时间平分为两部分。\par

4. \BibleQuote{他们作为\Emph{记号，并作为节令、日子和年岁}。}\BibleRef{创 1:14} 发光体所给的记号是人类生活所必需的。事实上，如果我们不带着过分的好奇心询问，长久的经验会使我们发现多么有用的观察！什么记号预示下雨、干旱、风的升起（局部的或普遍的，强烈的或温和的）！主在说到以下话时向我们指明了太阳所给的其中一个记号：\BibleQuote{今天必是坏天气，因为天色发红而阴暗。}\BibleRef{玛 16:3} 事实上，当太阳在雾中升起时，它的光线变得暗淡，但圆盘看起来像燃烧的炭，并呈血红色。正是空气的厚度造成了这种景象；由于太阳的光线不能驱散这样积聚而密集的空气，它肯定不能由从地面蒸发的水汽波所保留，并且将因过多的湿气在它所积聚的区域引发风暴。同样，当月亮被湿气包围，或当太阳被所谓的光环围绕时，这是大雨或暴风的预兆；同样，如果幻日伴随太阳运行，它们预告某些天体现象。最后，那些像彩虹颜色一样的直线出现在云上，预告下雨、异常的风暴，或者简而言之，天气的彻底变化。\par

那些致力观察这些天体的人，从月相的不同变化中找到了预兆，仿佛包裹大地的大气必须随着月相变化而改变。到了新月后的第三天，若月牙尖锐而清澈，便是持续晴朗的征兆。若其双角显得粗厚且泛红，则预示着要么有大雨，要么有南来的大风。谁不知道这些预兆在生活中有何等用处？多亏了它们，水手预见到风暴的危险，便将船只留在港湾；旅人则预先躲避祸害，等待天气好转。多亏了它们，忙于播种或栽培植物的农夫，能知道哪些季节有利于他们的劳作。此外，主也向我们宣告，在宇宙解体之时，太阳、月亮和星辰中将出现异兆：太阳将变为血，月亮不再发光，这些都是万物终结的征兆。\par

5. 然而，那些越过界限，以圣经的话为借口来推演命宫技艺的人，声称我们的生命取决于天体的运动，并说加色丁人因此能从行星中读出将要发生在我们身上的事。他们不理解这些极为简单的话\Quote{作为记号}，既不是指天气变化，也不是指季节更替；他们只根据自己的想像，在其中看到人类命运的分配。他们实际上说什么？当行星在黄道十二宫中交叉时，它们相遇所形成的某些图形产生某些命运，而其他图形则产生不同的命运。\par

也许为了清晰起见，详细讨论一下这种虚妄的学问并非无益。我不以自己的话来反驳他们；我将用他们的话，为受感染的人提供医治，为其他人提供预防坠落的保护。占星术的发明者看到在时间的广度中许多征兆逃过了他们，便将时间划分，将每一部分封闭在狭窄的界限内，仿佛在最微小、最短暂的间隔中，在眨眼之间——用宗徒的话说\BibleRef{格前 15:52}——一个出生与另一个出生之间竟有最大的差异。某人在这一刻出生：他将成为城市的王侯，治理人民，富足而有权势。另一人在下一刻出生：他将贫穷、悲惨，每日从这家到那家乞讨面包。因此，他们将黄道分为十二部分，并因太阳需三十天穿过这无误圆圈中的每个十二分之一，他们再将每个部分分为三十份。每一份再分为六十个新份，这些新份又再分为六十份。那么，试看，在确定一个婴孩的出生时，是否可能遵守这种严格的时间划分？孩子出生了。接生婆确认性别；然后她等待哭声，那是生命的标志。在这之前，你认为过了多少时刻？接生婆向加色丁人宣布孩子的出生：在她开口之前，你能数出多少分钟？尤其是如果记录时辰的人是在产房之外？而我们知道，无论白天黑夜，查询日晷的人都应以最精确的准确度标记时辰。在这段时间内，多少秒一掠而过！因为命宫的行星不仅应当落在黄道十二分之一中，甚至落在其第一个细分中，而且还要落在再分此细分的六十分之一中，并且，为了得到准确真理，还要落在它又包含的六十分之一细分中。而要获得如此精微的知识——从此刻起便不可能捕捉——必须询问每一颗行星，以确定其在黄道十二宫中的位置，以及行星在孩子出生时刻所形成的图形。因此，若不可能准确找到出生的时辰，且最微小的变化便能颠覆一切，那么，那些投身于这种想像出来的科学的人，以及那些张口倾听他们、仿佛能从中得知未来的人，都是极为可笑的。\par

6. 然而，产生了什么效果？某人生就卷发和明亮的眼睛，因为他生于公羊之下；公羊的外貌便是如此。他会有高尚的情感，因为公羊生来就统治。他会慷慨大方并富有资源，因为这动物毫不费力地脱去羊毛，大自然立刻又给它披上新的。另一人生于公牛之下：他将吃苦耐劳，具奴仆性格，因为公牛屈于轭下。另一人生于天蝎之下：像这有毒的爬虫，他将是个攻击者。生于天秤之下的人将是公正的，因为我们的天秤是公正的。这岂非愚妄之极？这公羊——你从中引出人的命宫——是天的十二分之一，太阳进入它时便到达春季。天秤和公牛也是黄道的十二分之一。你如何能在那里看到影响人生命的主要缘由？并且，你为何以动物来刻画进入这世界之人的品性？生于公羊之下的人将慷慨大方，不是因为天的这一部分给予这种特性，而是因为野兽的本性如此。那么，我们为何要因这些星辰的名字而自吓，并劝说自己相信这些咩咩叫声？如果天从这些动物获得不同的特性，那么它自己便受外在影响，因为它的原因取决于在田野中吃草的畜牲。这是可笑的论断；但更可笑的是，竟假装从毫无关联的事物中得出互相影响的结果！这所谓的科学是真正的蜘蛛网：若一只蚊虫或苍蝇，或某种同样微弱的昆虫落入其中，便被缠住；若有一只较强的动物靠近，它便穿过无碍，将薄弱的网一并带走。\par

7. 他们并不在此止步；甚而我们的行为——每个人都感觉到自己的意志在其中主宰，我指的是德行的实践或恶行的实践——按照他们的说法，也取决于天体的影响。认真反驳这种错误本是可笑的，但由于它网罗了许多人，也许最好还是不要默不作声。我首先要问他们：星辰描绘的形状难道不是一天变化上千次吗？在行星的永恒运动中，有的运行更快，有的运转较慢，常常在一小时内我们看见它们相互对视然后又隐藏起来。在出生的时辰，是否被一颗善星或一颗恶星注视——用他们的语言来说——至关重要。占星家们常常抓不住善星显现的时刻，由于让这瞬息的机会溜走，他们就将新生儿登记在恶星的影响下。我不得不使用他们自己的话。何等疯狂！但尤其是何等不敬！因为恶星将自己的邪恶归咎于造它们的祂。如果邪恶内在于它们的本性，那么造物主就是邪恶的制造者。如果它们自己制造邪恶，那它们就是被赋予选择能力的动物，其行为将是自由且自愿的。对无灵魂的存在说这些谎言，难道不是愚蠢之极吗？再者，不分个人功过而分配善恶，说一颗星因其所在位置而成为善星，因被另一颗星注视而变成恶星，如果它从这图形稍稍移动就失去其恶的影响——这显示出多大的缺乏理性呢！\par

但让我们继续。如果在时间的每个瞬间，星辰都变化其形状，那么在一天众多时刻的这千变万化中，应该会重现王室出生的配置。那么为什么不是每天都有国王诞生？为什么王位是从父亲到儿子继承？毫无疑问，从来没有一位国王采取措施让他的儿子在星宿下出生为君王。因为有什么人拥有这样的权力？那么\Person{乌齐雅}如何生了\Person{约堂}，\Person{约堂}生了\Person{阿哈兹}，\Person{阿哈兹}生了\Person{希则克雅}？他们中又没有一个人出生在奴隶的时刻，这是出于什么偶然？如果我们德行与恶行的根源不在我们自己，而是我们出生的宿命结果，那么立法者为我们规定当做和不当做的事就是无用的；法官尊崇德行而惩罚恶行也是无用的。强盗没有罪过，刺客也没有罪过：那是为他所注定的；他不可能缩回他的手，因为被不可避免的必然性催迫去行恶。那些辛苦致力技艺的人是最疯狂的人。农夫无需播种、无需磨镰就能获得丰收。商人无论愿意与否都会发财，并会被命运淹没在财富中。至于我们基督徒，我们将看到我们伟大的希望破灭，因为人若不自由行动，就没有公义的赏报，也没有罪恶的惩罚。在必然与命运的统治下，没有功绩的位置，而功绩是一切公正审判的首要条件。但我们且住。你们自身健全的人无需再听，时间也不允许我们无限期地攻击这些不幸的人。\par

8. 让我们回到后面的话。\BibleQuote{作标记，定节令、日子和年。} \BibleRef{创 1:14} 关于标记我们已经谈过。至于时节，我们理解为季节的交替——冬、春、夏、秋，我们看到它们如此规律地相互接续，这得益于星辰运动的规律性。当太阳南方逗留，给我们的地区带来大量夜影时，便是冬天。弥漫在地面的空气寒冷，聚集在我们头顶的湿汽产生雨水、霜冻和无数雪花。当太阳从南方返回，行至中天，将昼夜平分，它越长时间停留在地面之上，就越为我们带回温和的气温。于是春天来临，它使所有植物发芽，使大部分树木恢复生机，并通过连续繁殖延续所有陆生和水生动物。此后，太阳向北回归线返回，给我们最长的白昼。当它在空中走得更远时，它烘烤我们头顶之上，使地面干燥，使谷物成熟，加速树木果实的成熟。在最热的时期，太阳在正午投下的阴影很短，因为它从我们头顶上方照射。所以最长的白昼就是阴影最短的时候，同样，最短的白昼就是阴影最长的时候。这正是我们所有居住在地球北部地区的\OriginalTerm{偏影}{Hetero-skii}（即单侧有影的人）所遇到的。但有些人一年中有两天在正午完全没有阴影，因为太阳垂直在他们的头顶，从四面均匀地照亮，以至于它可以通过窄口照到井底。因此有人称他们为\OriginalTerm{无影}{askii}（即无影的人）。至于那些生活在香料之地之外的人，他们的影子有时偏这边，有时偏那边，是这片土地上唯一在正午阴影落在这边或那边的人；因此他们被称为\OriginalTerm{双侧影}{amphiskii}（即双侧有影的人）。所有这些现象都发生在太阳进入北方地区时：它们让我们感受到太阳光线投射到空气中的热量及其产生的效果。接着我们过渡到秋天，它打破过热，逐渐减少温暖，并以适中的气温毫无痛苦地将我们带回冬天，即太阳从北方返回南方的时候。这样，季节跟随太阳的路径交替出现，来支配我们的生活。\par

\Quote{它们要成为日子} \BibleRef{创 1:14}，圣经说，不是要产生它们，而是要管理它们；因为白天和黑夜比光体的受造更古老，这正是圣咏向我们宣告的。\Quote{太阳管理白天……月亮和星辰管理黑夜。}太阳如何管理白天？因为它到处携带光亮，一旦升上地平线，便驱散黑暗，给我们带来白天。这样，我们可以毫不自欺地定义白天为被太阳照亮的大气，或太阳在我们半球经过的时间。太阳和月亮的功用更进一步用以标记年份。月亮在运行十二次之后，形成一年，有时需要一个闰月以使其与季节完全一致。这就是从前希伯来人和早期希腊人的年。至于太阳年，则是太阳从某个星宫出发，按正常进程回到该星宫所需的时间。\par

9. \Quote{\Emph{天主造了两个大光体}} \BibleRef{创 1:16}。\Quote{大}这个词，例如当我们说天空、大地或海洋的伟大时，可能具有绝对意义；但通常它只有相对意义，如大马或大牛。这并不是说这些动物体型过大，而是与同类相比，它们配得上称为大。我们自己在这里对伟大应形成何种概念？是否该是我们在蚂蚁和自然中一切小生物身上所认识的那种概念，我们称它们为大是与同类相比，并显示它们超越同类？或者我们应将伟大归于光体，如同归于它们内在的自然伟大？至于我，我这样认为。若太阳和月亮是大的，并非与较小的星辰相比，而是因为它们的尺寸如此之大，以致它们所散发的光辉照亮了天穹和大气，同时拥抱了大地和海洋。无论它们在天上的哪个部分，或升起，或落下，或在天顶，它们在人眼中总是显得一样大，这是它们巨大尺寸的明显证明。因为整个天穹不能使它们在某一处显得更大，而在另一处显得更小。我们远看的对象在眼中显得矮小，而当它们接近我们时，我们就能对它们的大小形成更准确的概念。但没有人能离太阳更近或更远。大地上的所有居民从同样的距离看它。印度人和不列颠人看它同样大小。东方人看它落下时不觉得它变小；西方人看它升起时不觉得它更小。若它在天正中，则两方面都无变化。不要被表面现象所欺骗，因为它看上去有一肘宽，就以为它没有更大。在非常远的距离，物体总在我们眼中失去大小；视觉不能越过中间空间，犹如在半途中耗尽，只有一小部分到达可见对象。我们的视力很小，使得我们看到的一切都显得很小，以自己的状态影响所见之物。因此，若视觉出错，它的见证就不可靠。回想你自己的印象，你会在自己身上找到我话语的证据。若你曾从高山顶上俯瞰一片宽阔平坦的平原，牛群在你看来有多大？农夫本身有多大？它们岂不看来如蚂蚁？若从巍峨的岩石顶上，俯视广阔的大海，你放眼望去，最大的岛屿在你看来有多大？那些大吨位的船只，在蓝色海面上展开白帆，在你看来有多大？它岂不比鸽子还小？这是因为视觉，正如我刚才告诉你的，在空气中迷失，变得微弱，不能精确把握所见之物。而且进一步：你的视觉向你显示高山被山谷切割得圆润平滑，因为它只达到突出部分，由于软弱，不能穿透分隔它们的山谷。它甚至不能保持对象的形状，以为所有方塔都是圆的。因此，这一切都证明在远处视觉只给我们呈现模糊混乱的对象。那么，按圣经的见证，光体是大的，且远比其外表无限更大。\par

10. 再看另一个它伟大的明显证明。虽然天穹可能充满无数星辰，但它们所有的光都不能驱散黑夜的阴暗。太阳独自，从它出现在地平线的那一刻起，当它还在被期待、尚未完全升上大地之上时，便驱散了黑暗，胜过众星，消散并扩散了此前在我们头上厚重而凝聚的空气，并由此产生了晨风和晴朗天气中漫过大地的露水。难道大地如此广阔，若不是一个巨大的圆盘将光芒倾泻其上，怎能在一瞬间完全被照亮？在此请认出造物主的智慧。看祂如何使太阳的热量与这一距离相称。它的热量如此调节，既不因过度而灼烧大地，也不因不足而使大地变冷而贫瘠。\par

凡此种种，皆与月亮的特性相类：月亮亦有一个巨大的身体，其光辉仅逊于太阳。然而，我们的眼睛并不总能看见她的全貌。时而她呈完美的圆盘状，时而又因缺损而显得一侧不足。盈时一侧蒙影，亏时另一侧隐没。万有的神圣造物主并非无端使月亮在不同时期呈现如此不同的形态。这为我们揭示了人性的一个鲜明例证：人身上没有什么是稳固的；从虚无中，他升向完美；然而，一旦急于施展力量达至其全盛，便骤然陷入逐渐衰退，因消减而消亡。因此，看见月亮，使我们想到人事的迅疾变迁，理应教导我们不要因今世之善而自傲，不要炫耀自己的权势，不要被不可靠的财富冲昏头脑，要轻视那易变的肉体，而照料灵魂，因为灵魂的善是固定不移的。你若不能毫无悲伤地看待月亮在不知不觉中逐渐亏损失去光辉，那么，当看见一个灵魂，曾拥有德行，却因疏忽而失去美丽，不恒久于其情感，反倒因志向不定而动荡不安、不断变化时，你岂不更应伤痛？圣经所言极其真实：\BibleQuote{愚人如月亮般变化。}\BibleRef{德 27:11}\par

我也相信，月亮的盈亏并非不对动物及一切生物的机体产生巨大影响。这是因为身体在其盈与亏时状态不同。当月亮亏损时，身体失去密度而变得虚空；当月亮渐盈且接近满月时，身体似乎同时充盈，这是由于月亮散发出一种渗透一切的、混合着热气的细微湿气。证据如下：看那些在月光下睡眠的人，头部充满丰沛的湿气；看鲜肉在月亮作用下如何迅速变质；看动物的脑、海洋生物最潮湿的部分、树木的髓。显然，月亮必如圣经所言，具有极大的体积与能力，才能使整个自然如此参与她的变化。\par

11. 月亮的盈亏也决定了空气的状况，这由新月过后常出现的突然扰动所证明——在风平浪静之时，云层被搅动并相互碰撞；海峡中的潮汐涨落以及海洋的潮汐亦为明证，以致沿海居民看到海洋规律地跟随月亮运转。海峡之水在月亮的各个相位期间，从一侧涌向另一侧，又从另一侧退去；但当月亮是新月时，海水不得片刻安宁，在永久的往返摇摆中运动，直到月亮重现，将潮汐调节有序。至于西方的海洋，我们看到其潮汐时而退回海床，时而漫溢，这是月亮以其呼气将其吸回，再由呼气将其推向自己的界限。\par

我详述这些细节，是为向你们展示日月星辰的宏伟，并使你们看到，在受感的话语中没有一句是虚空的。然而，我的讲道几乎尚未触及任何重要之点；关于太阳和月亮的大小与距离，还有许多其他发现，任何人若认真研究它们的作用与特性，均可借助理性获得。因此，请容许我坦诚承认自己的软弱，以免你们以我的话语来衡量造物主的大能工作。我所言及的这一点点，反而应使你们推测出我省略未述的奇妙之处。因此，我们不应以眼睛，而应以理性来衡量月亮。理性在发现真理方面，远比眼睛可靠。\par

到处流传着荒唐的老妪妄谈，想象于醉酒的谵妄之中，诸如：魔法能使月亮离开原位，降到地上。巫师的法术怎能使至高者所立根基的月亮动摇？若它一旦被拔除，有何处能容纳它呢？\par

你们是否想从微小的迹象得到月亮大小的证据？世上所有的城镇，无论相距多么遥远，都同样从月亮那里获得光照，照射在那些朝向月亮升起的街道上。如果月亮不是面向所有人，那么只有正对着她的地方才会完全照亮；至于在她圆盘两端之外的地方，则只能接收到偏离和偏斜的光线。这就是屋内灯光产生的效果：若一盏灯周围有数人，只有直接面对灯的人投下笔直的阴影，其余的人则沿两侧倾斜的线投下阴影。同样，如果月亮的天体不是巨大无比，她就不可能平等地照临一切。实际上，当月亮在赤道区域升起时，所有的人都平等地享受她的光，无论他们是居住在北极附近冰带之下的人们，还是居住在极南靠近热带地区的人们。月亮通过似乎与所有人面面相对的方式，向我们展示了她的巨大。那么，谁能否认一个如此广大而能平等分施其光的天体之宏大呢？\par

关于太阳和月亮的伟大，已经说得够多了。愿那赐予我们智慧，使我们在受造的最小物体中认出\OriginalTerm{造物主}{Contriver}之大智慧的天主，使我们在庞大的物体中，对它们的创造者产生更高的认识。然而，与它们的创造者相比，太阳和月亮不过是苍蝇和蚂蚁。整个宇宙不能给我们关于天主伟大的正确认识；我们只是借着自身微弱而细小的标志，往往借助于最小的昆虫和最微不足道的植物，才能提升自己归向祂。满足于这些话，让我们献上感谢：我感谢那赐予我\OriginalTerm{圣言的职务}{ministry of the Word}的祂，你们感谢那以神粮喂养你们的祂；祂甚至在此刻，使你们在我微弱的声音中找到了大麦饼的力量。愿祂永远喂养你们，并按照你们的信德，赐给你们\OriginalTerm{圣神的彰显}{manifestation of the Spirit}在\Person{耶稣基督}、我们的主内。愿光荣与权能归于祂，至于无穷之世。阿们。\par

\SourceRecord{Translated by Blomfield Jackson.；Nicene and Post-Nicene Fathers, Second Series；Vol. 8.；Edited by Philip Schaff and Henry Wace.；Buffalo, NY: Christian Literature Publishing Co.,；1895.；<http://www.newadvent.org/fathers/32016.htm>.}

% structure: p0001 source=h1 target=section reason=page-title

\section{六日创造（第七讲）}

\Emph{蠕动的生物的创造}\par

1. \BibleQuote{\Emph{天主说：水中要繁生蠕动的生物}}，各按其类，\BibleQuote{\Emph{并有飞鸟飞翔在地面上}}，\Emph{各按其类}。\BibleRef{创 1:20} 在光体被造之后，如今水中充满了生物，世界这一部分也得到了自己的装饰。大地从自己的植物中获得了装饰，诸天得到了星辰的花朵，大光体如同两只眼睛，协调地美化它们。还剩下水要接受装饰。命令一下，江河湖泊立即成为肥沃，产生了它们天生的幼崽；海洋孕育了各种游泳的生物；甚至泥沼和水洼中的水也没有闲置；它参与了创造。到处从水的沸腾中出现了青蛙、蚊虫和苍蝇。因为今天我们所见的，是过去的标记。这样，水到处都急于服从造物主的命令。谁能数算天主伟大而不可言喻的能力所使忽然出现的有生命的、活动着的物种呢？当这个命令授权水产生生命时，让水产生蠕动的生物。这是第一次造出有生命和感觉的存在。因为虽然植物和树木被称为有生命的，因为它们分有被滋养和生长的能力，但既不是有生命的生物，也没有生命。为了创造这些，天主说：\Quote{让水产生蠕动的生物。}\par

每一个会游泳的生物，无论它在水面上滑行，还是劈开深处，都具有活动物的性质，因为它在水体上拖行。某些水生动物有脚并能行走；尤其是两栖动物，如海豹、螃蟹、鳄鱼、河马和青蛙；但它们以上游泳的能力为首要。因此说：\BibleQuote{让水产生蠕动的生物。}在这几句话中，遗漏了哪个物种？哪个没有被包括在造物主的命令中？难道我们没有看见胎生动物、海豹、海豚、鳐鱼和所有软骨鱼吗？难道我们没有看见卵生动物，包括各种鱼类，有皮的和有鳞的，有鳍的和无鳍的吗？这个命令只需要一句话，甚至不到一句话，一个记号，一个天主意愿的示意，便具有如此宽广的含义，以至于包含了所有鱼类的各种类和科。要一一列举它们，就是企图数算海洋的波浪或用手心测量它的水。\BibleQuote{让水产生蠕动的生物。}也就是说，那些充满深海的和那些喜爱岸边的；那些居住在深处的和那些附着在岩石上的；那些群居的和分散生活的，鲸类、巨大的和微小的。都从同一能力、同一命令中，大大小小都获得存在。\BibleQuote{让水产生蠕动的生物。}这些话向你显示了水中游泳动物的自然亲和性；因此，鱼一出水就很快死亡，因为它们没有呼吸能吸引我们的空气，水是它们的元素，如同空气是陆地动物的元素一样。原因很清楚。在我们身上，肺，那多孔且有海棉状的内脏部分，通过胸部的扩张接受空气，分散并冷却内部的热量；在鱼身上，鳃的运动，通过交替开合来吸入和排出水，代替了呼吸。鱼有特殊的命运、特殊的本性、自己的营养、独立的生活。因此它们不能被驯服，也不能忍受人的手触摸。\par

2. \BibleQuote{让水产生蠕动的生物，各按其类。} 天主使每个物种的首生者出生，作为自然的种子。它们数量众多的数目在随后的世代中得以保持，当它们需要生长和繁衍时。另一种类是甲壳类，如贻贝、扇贝、海螺、贝螺，以及无穷无尽的各种牡蛎。另一种类是甲壳纲，如螃蟹和龙虾；另一种是软体无壳的鱼类，肉质柔软细嫩，如章鱼和乌贼。而在这些之中，多么不可胜数的种类！有鲂鱼、七鳃鳗和鳗鱼，产生于河流和池塘的泥中，其本性更像有毒的爬行动物而不是鱼。另一种类是卵生类；另一种是胎生类。后者包括剑鱼、鳕鱼，总之，所有软骨鱼，甚至大部分鲸类，如海豚、海豹，据说如果它们看到自己的幼崽还很小，受到惊吓，就把它们吞回腹中保护它们。\par

\BibleQuote{\Emph{让水各按其类产生。}} 鲸类是一个物种，小鱼是另一个物种。不同种类中有何等无穷的多样性！它们都有自己的名字，不同的食物，不同的形状、外表和肉质。所有都呈现无穷的多样性，并被分为无数的类别。有没有一个金枪鱼渔夫能向我们列举那种鱼的不同品种？然而他们告诉我们，看到大群鱼时，他们几乎能说出组成鱼群的个体数目。在所有那些在海岸和岸边度过漫长一生的人中，有谁能确切地告诉我们所有鱼类的历史呢？\par

有些鱼类为\Place{印度洋}的渔夫所熟知，另一些则为\Place{埃及海湾}的劳动者所知晓，还有一些为岛民所认识，另外一些则为\Place{毛里塔尼亚}人所了解。大小鱼类都因这最初的命令、因这不可言喻的能力而被同样创造。它们的食物何其不同！每个物种繁殖的方式何其多样！大多数鱼不像鸟那样孵卵；它们不筑巢，也不辛苦喂养幼鱼；是水接收并在其中激发卵的生命。对于它们而言，每个物种的繁殖是不变的，本性不会有混杂。没有地上那种产生骡子和某些违反其物种本性的鸟类的结合。鱼类中，没有像牛和羊那样配备半副牙齿的品种，除了某些作者所说的scar鱼外，没有反刍的。所有鱼都有紧密而非常锋利的牙齿，唯恐食物因咀嚼时间过长而逃脱。事实上，如果食物不被立即咬碎和吞下，就会被水冲走。\par

3. 鱼类的食物依其种类而异。有些以泥为食；有些吃海草；另一些满足于水中生长的草。但大部分互相吞食，较小的成为较大的食物。若一条鱼捕获了比自己更弱的鱼，却又成为另一条鱼的猎物，那么征服者和被征服者最终都葬身于后者的腹中。我们凡人，在压迫下级时，不也如此行事吗？那最后一条鱼和那个被贪得无厌的欲望所驱使、用其无法满足的贪婪之网吞没弱者的人，有什么区别？你曾占有穷人的财物；你捕获了他，并使他成为你丰裕的一部分。你比不义之人更不义，比吝啬之人更吝啬。当心，免得你像鱼一样，最终被钩、被笼、或被网捕获。的确，当我们做了恶事之后，也终究难逃惩罚。\par

现在请看，一个弱小的动物身上有多少诡计和狡猾，并学会不要效法作恶者。螃蟹爱吃牡蛎的肉；但牡蛎有壳作为庇护，这是大自然赋予其柔软细腻肉质的坚固壁垒，使其难以捕获。因此他们称牡蛎为\OriginalTerm{碎壳皮}{sherd-hide}。由于它被两片完美契合的壳包裹，螃蟹的钳子完全无能为力。它怎么办呢？当它看到牡蛎避开风，愉快地取暖，半张开壳向着太阳时，它偷偷地投入一颗小石子，阻止壳合拢，从而用诡计取得了用武力无法得到的东西。这些缺乏理性和言语的动物，其恶意竟至于此。但我希望你能在狡猾和巧诈上与螃蟹一争高下，同时避免伤害邻人；这只螃蟹正是那狡猾地接近兄弟、利用邻人的不幸、以他人痛苦为乐之人的形象。啊，不要效法那被定罪者！满足于你自己的命运。在智慧者眼中，拥有必需品的贫穷胜于一切享乐。\par

我不会缄默不言乌贼的狡诈和诡计，它会依附岩石并呈现岩石的颜色。大多数鱼漫不经心地游向乌贼，如同游向岩石，结果自己成了这奸猾之物的猎物。这样的人便是那些阿谀权贵之人，他们随时随势而变，不在一瞬间持守同一目的；他们在节制者面前称赞节制，在放纵者面前称赞放纵，迎合每个人的喜好。要逃脱他们并对他们的恶行保持警惕是困难的；因为他们在友谊的面具下隐藏其精明的邪恶。这样的人是披着羊皮的贪婪之狼，正如主所称他们的那样：\BibleQuote{披着羊皮的贪婪之狼}\BibleRef{玛 7:15}。要逃避摇摆不定和见风使舵；要追求真理、真诚和纯朴。蛇是狡猾的；所以它被诅咒用肚皮爬行。义人是正直的人，如同约伯。因此天主使孤独的人有家居住\BibleRef{咏 68:7}。这就是那广大无边的海，其中有无数爬行的生物，大小兽类都有。然而，这些动物中存在着智慧和奇妙的秩序。鱼类并不总该受我们的责备；它们常常给我们有益的榜样。为什么每一种鱼都满足于被分配的区域，从不越过自己的界限进入外海？没有测量员为它们分配住处，也没有用墙围起它们，或为它们划定界限；每个种类都被自然地分配了家园。一个海湾养育一种鱼，另一个海湾养育其他种类；那些在此处繁盛的种类在别处却不存在。山岳没有在它们之间耸起尖峰；河流也没有阻挡它们的去路；这是自然法则，它根据每个种类的需要，以平等和公正为它们分配了居所。\par

4. 我们却并非如此。为什么？因为我们不断移动我们祖先所立定的古时地界\BibleRef{箴 22:28}。我们侵占，我们加添房屋于房屋，田地于田地，靠损害邻人自肥。大鱼知道大自然给它们指定的居所；它们占据远离人类居所的海域，那里没有岛屿，也没有大陆阻挡它们，因为从未有人航行过，好奇心和需要都没有说服水手去冒险。居住在这片海中的巨兽，其大小如同高山，据目击者所说，它们从不越界去蹂躏岛屿和沿海城镇。这样，每个种类仿佛都驻扎在城镇、村庄、古老的土地上，其居所就是分配给它们的海域区域。\par

然而，人们也曾见过洄游鱼类的实例，它们仿佛经过共同商议，被带往异域，在某个信号下全部启程。当繁殖的季节到来时，它们仿佛被自然的共同法则唤醒，从一个海湾迁移到另一个海湾，朝着\Place{北海}方向前进。在它们返回的时候，你可以看见所有这些鱼如同洪流一般穿过\Place{普洛庞提斯海}，涌向\Place{攸克辛海}。是谁将它们列队布阵？哪里有君王的命令？难道是公告牌上张贴的谕令指示了它们出发的日子？谁为它们充当向导？请看天主的命令如何涵盖一切，并延伸至最微小的对象。一条鱼不反抗天主的律法，而我们人类却不能忍受祂救恩的诫命！不要因鱼是哑巴且毫无理性而轻视它们；反倒要惧怕，在你反抗造物主安排的时候，你比它们更缺乏理性。听听这些鱼，它们以行动几乎在说话，说：是为了我们种族的延续，我们才进行这次长途航行。它们没有理性的恩赐，却有着坚定内在的自然律，指示它们该做什么。\Quote{让我们去吧，它们说，到\Place{北海}去。}那里的水比其他海域更甜；因为太阳在那里停留不长，其光线没有吸走所有可饮用的部分。即便是海洋生物也喜爱淡水。因此，人们常见它们进入河流，从海里向上游很远。这就是为什么它们偏爱\Place{攸克辛海}胜过其他海湾，因为那里最适合繁殖和养育幼鱼。当它们达成目的后，整个族群便返回家园。让我们听听这些哑巴生物告诉我们原因。\Place{北海}，它们说，水浅，表面暴露在强风中，海岸和避风处很少。因此，风很容易搅动其海底，将沙粒与波浪混合。此外，冬季寒冷，因四面八方有大量河流注入。因此，在夏季适度享用其水域之后，当冬季来临时，它们便急忙前往较暖的、被阳光加热的深海处，逃离北方的风暴区域，在较平静的海域中寻求港湾。\par

5. 我亲眼见过这些奇迹，并惊叹于天主在万有中的智慧。如果无理性的受造物能够思考并为自己的保存做准备；如果一条鱼知道它应当追求什么、避开什么，那么，我们这些享有理性、受法律教导、被应许鼓励、因圣神而明智的人，却为何在关乎自己的事上比鱼更缺乏理性？它们懂得为未来筹划，我们却放弃对未来的希望，在兽性的放纵中消耗一生。一条鱼穿越整个海洋寻找对它有益的东西；那么，你——生活在万恶之源的闲散中——将怎样说呢？不要让任何人以无知为借口。在我们内已植入了自然理性，它告诉我们与善认同，并避开一切有害的。我不必远离大海就能找到例子，因为那正是我们研究的对象。我曾听一位住在海边的人说，海胆——一种微不足道的可鄙生物——常常向水手预报风平浪静或暴风骤雨。当它预感到风的骚动时，便躲到一块大卵石下，紧紧附着其上如同锚，安然翻滚，因重量而稳住，不致成为波浪的玩物。这对水手来说是确凿的征兆，表明他们面临狂风激荡的危险。没有一位占星家或迦勒底人，通过星辰升起解读空气的扰动，曾将自己的秘密传授给海胆：是海洋与风的主宰将这伟大智慧的明显证明印在这小动物身上。天主预知一切，祂没有忽略任何事物。祂永不闭目的眼注视着万有。祂处处都在，赐予每个受造物保存自身的方法。如果天主没有将海胆置于祂的眷顾之外，祂会不关心你吗？\par

\BibleQuote{\Emph{丈夫们，要爱你们的妻子}}。\BibleRef{弗 5:25} 虽然你们由两个身体组成，但被结合，生活在婚姻的共融中。愿这自然的纽带，愿这因祝福而加于你们的轭，使分离者重新合一。毒蛇，最凶残的爬行类，与海鳝结合，并以嘶声宣告自己的存在，从深处召唤它前来进行夫妻结合。海鳝服从，并与这有毒的动物结合。这意味着什么？无论丈夫多么难相处、多么凶暴，妻子都应当忍耐他，不愿寻找任何借口破坏结合。他击打你，但他是你的丈夫。他是个醉汉，但他与你因本性而结合。他粗暴易怒，但从今以后他是你肢体的一部分，且是最珍贵的。\par

6. 让丈夫们也听听：这里有一个教训给他们。毒蛇出于对婚姻的尊重，吐出毒液；而你，难道不会出于对你结合体的尊重，除去你灵魂中的残暴与非人性吗？或许毒蛇的例子还有另一层含义。毒蛇与海鳝的结合是对自然的奸淫性侵犯。你，这个在暗中破坏他人婚姻的人，要知道你像什么爬虫。我只有一个目的，就是使我所说的一切都服务于教会的建立。那么，让放荡之人约束他们的情欲，因为他们被地上和海中生物的榜样所教导。\par

我的身体虚弱和天色已晚迫使我结束我的讲道。然而，关于海中的产物，我仍有许多观察要向专心听讲的你们表达，以激发你们的赞叹。谈论海本身：海水如何变成盐？珊瑚这种备受珍视的石头，在海中本是一种植物，一旦暴露于空气便坚硬如石？为何大自然将如此珍贵的珍珠包裹在最卑微的动物——牡蛎之中？这些被帝王宝匣垂涎的珍珠，被抛掷于海岸、礁石上，由牡蛎壳包裹。海蚌如何产出无人能仿制的金色绒毛？贝壳如何赐给帝王紫红色，其光彩不亚于田野的花朵？\par

\Quote{\Emph{让水滋生}}。有什么必需之物没有立即出现？有什么奢侈品没有赐予人？有些供应人的需要，有些使人默观创造的奇迹。有些令人恐惧，以教训我们的懒惰。\BibleQuote{天主创造了巨鲸。}\BibleRef{创 1:21}圣经称它们为\Quote{巨大的}，并非因为它们比虾和鳀鱼更大，而是因为它们的身体大小如同大山。因此，当它们在水面游动时，人们常看到它们像岛屿一样出现。但这些巨兽并不出没于我们的海岸，而是栖息在大西洋。这些动物被创造出来，是为了使我们惊恐和敬畏。若你现在听说，最大的船只张满风帆，却被一种极小的鱼——吸盘鱼轻易阻止，以致船只在海中长久不动，仿佛生了根，难道你不在这小东西身上看到造物主大能的同样证明吗？因此，令人畏惧的不仅是剑鱼、锯鳐、角鲨、鲸和鲨鱼；我们同样害怕黄貂鱼死后的尖刺，以及海兔那迅速而不可避免的致命一击。如此，造物主愿意这一切都使你保持警醒，使你满怀对祂的希望，逃避这些生物威胁你的灾祸。\par

但让我们离开海的深处，到岸边避难。因为创造的奇迹如波浪般接连不断，淹没了我的讲论。然而，我若在陆地上发现更伟大的奇迹，我的灵魂或许会像约纳一样逃往海中。但我似乎感到，遇见这些数不清的奇迹使我忘记了一切尺度，如同那些在汪洋中航行而没有固定航标的人，常常不知自己航行了多远。这正是我所经历的：当我的话语掠过创造时，我未曾觉察到向你们讲述的受造物如此众多。尽管这位尊贵的会众乐于听我的讲论，并且叙述主人的奇迹使仆人耳悦，但容我将讲论的船在此抛锚，等待明天向你们交付其余部分。因此，让我们全都起身，为已听讲的感谢，并祈求力量以聆听其余。在用餐时，愿你们餐桌上的谈话围绕今早和今晚所论之事。充满这些思想，即使在睡眠中，你们也能享受白日的愉悦，以便能说：\BibleQuote{我睡眠，但我的心却醒着，}\BibleRef{歌 5:2}日夜默思主的法律，愿光荣与权能归于祂，至于无穷之世。阿们。\par

\SourceRecord{Translated by Blomfield Jackson.；Nicene and Post-Nicene Fathers, Second Series；Vol. 8.；Edited by Philip Schaff and Henry Wace.；Buffalo, NY: Christian Literature Publishing Co.,；1895.；<http://www.newadvent.org/fathers/32017.htm>.}

% structure: p0001 source=h1 target=section reason=page-title

\section{\WorkTitle{六日创世（第八篇讲道）}}

\Emph{飞禽与水生动物的创造。}\par

1. 天主说：\BibleQuote{\Emph{地上要生出活物，各按其类：牲畜、爬虫和地上的野兽，各按其类；事就这样成了。}} \BibleRef{创 1:24} 天主的命令逐步推进，大地因此接受了它的装饰。昨天曾说：\Quote{让水产生爬行之物}，今天却说：\Quote{让地上生出活物。} 难道大地是有生命的吗？那些头脑疯狂的摩尼教人给大地赋以灵魂，难道是对的吗？听到这些话：\Quote{让地上生出活物}，大地并非产生其中所含的种子，而是那发令者同时赐予它生出活物的恩宠和能力。当地听到这话：\Quote{让地上生出青草和结果实的树木}时，并非它隐藏的草使它发芽，它也没有将一直藏在其深处的棕榈树、橡树、柏树带到表面。是天主的话语形成了受造物的本性。\Quote{让地上生出}，意思是说：不是让它生出它所拥有的，而是让它获得它所缺少的，当天主赐给它能力时。同样，现在：\Quote{让地上生出活物}，不是产生大地本身所包含的活物，而是天主的命令所赐予的活物。此外，摩尼教人自相矛盾，因为如果大地产出了生命，它就使自己失去了生命。他们可憎的教义无需证明。\par

但为什么水接受命令产生有生命的爬行之物，而地上产生活物呢？我们推论，按本性，水栖动物似乎只有不完美的生命，因为它们生活在厚厚的水元素中。它们听觉迟钝，视力模糊，因为透过水看东西；它们没有记忆、想象、社交意识。因此，神圣的语言似乎表明，在水生动物中，肉身生命产生了灵魂的活动；而在陆生动物中，由于具有更完美的生命，灵魂享有至高的权威。事实上，大部分四足动物的感官有更强的穿透力；它们对当前事物的感知敏锐，并精确地保持对过去的一切记忆。因此，似乎天主在命令水产生有生命的爬行之物后，为水生动物创造了简单的生命体，而为陆生动物，他命令灵魂存在并指导身体，表明地上的居民被赋予了更大的生命力。毫无疑问，陆生动物缺乏理智。但同时，它们每一个都通过自然的声音表达了灵魂的多少情感！它们用叫声表达快乐和悲伤，对熟悉事物的识别，对食物的需求，与同伴分离的遗憾，以及无数的情绪。相反，水生动物不仅无声，而且无法驯服、教导、训练它们与人类交往。\BibleQuote{牛认识自己的主人，驴也认识主人的槽。} \BibleRef{依 1:3} 但鱼不认识喂养它的人。驴认识熟悉的声音，认识常走的路，甚至人迷路时，它有时会充当向导。它的听觉比任何其他陆生动物更敏锐。哪种海中的动物能像骆驼那样表现出如此多的怨恨和愤怒？骆驼在被打后长时间隐藏怨恨，直到找到机会，然后报复。听啊，你们这些不宽恕的人，你们这些把复仇当作美德的人；看看你们像什么，当你们对邻居长时间保持愤怒，如同隐藏在灰烬中的火花，只等待燃料使你们的心燃烧起来！\par

2. \BibleQuote{\Emph{让大地生出活的生灵}}。为什么大地产出活的生灵？好让你能分辨牲畜的灵魂和人的灵魂。你很快就会知道人的灵魂是如何形成的；现在请听缺乏理性的受造物的灵魂。按照圣经，\BibleQuote{一切生物的生命在血内}，因为血在凝结时变成肉，肉在腐败时分解为土，所以野兽的灵魂本质上是土性的。\BibleQuote{让大地生出活的生灵}。看灵魂与血、血与肉、肉与土的亲缘关系；逆序而上，从土到肉，从肉到血，从血到灵魂，你会发现野兽的灵魂就是土。不要以为它比身体的本性更古老，也不要在肉分解后仍存留；避开那些傲慢哲学家的胡言，他们不羞于将自己的灵魂与狗相比；他们说自己从前曾是女人、灌木、鱼。他们曾否是鱼？我不知道；但我不怕肯定，在他们的著作中，他们的见识还不如鱼。\BibleQuote{让大地生出活的生物}。也许你们许多人问，在我快速演讲中为何有如此长的沉默？我听众中较好学的人不会不知道我词穷的原因。怎么！我没有看到他们互相使眼色，做手势让我看他们，提醒我遗漏了什么吗？我忘记了一部分受造物，且是最重要的一部分，我的讲论几乎结束而未曾触及它。\BibleQuote{让水多多滋生各种有生命的动物，和飞鸟，飞在地上，天空之中}。\BibleRef{创 1:20} 我在黄昏时讲了鱼；今天我们转到地上动物的考察；在这两者之间，鸟被我们遗漏了。我们像旅行者，忘记了一些重要物件，尽管已在路上很远，却不得不折回，因路途劳累而受疏忽的惩罚。所以我们不得不回头。我们所遗漏的不可轻视。它是动物受造中的第三部分，如果确实有陆地、有翼和水族三类动物的话。\par

\BibleQuote{\Emph{让水}}经上说：\BibleQuote{\Emph{多多滋生各种有生命的动物，和飞鸟，飞在地上，天空之中}}。为什么水也生出鸟类？因为可以说，飞鸟和水族之间有家族联系。正如鱼划开水面，用鳍向前推进，用尾鳍控制来回和直行运动，我们看见鸟借助翅膀在空中漂浮。两者都有游泳的特性，共同源于水使它们成为一族。同时，没有鸟是没有脚的，因为它所有的食物都在地上，不能缺少脚的役使。猛禽有尖爪，以便抓住猎物；其余的则被赋予了脚不可或缺的服务，以寻找食物并满足生活的其他需要。有少数鸟行走困难，它们的脚既不适合行走也不适合捕食。其中包括燕子，不能行走和寻找猎物，以及称为雨燕的鸟，它们以空气携带的小昆虫为生。至于燕子，其掠地的飞行履行了脚的功能。\par

3. 还有无数种类的鸟。如果我们像部分考察鱼那样全部回顾它们，就会发现，在同一个名字之下，飞行的受造物在大小、形状和颜色上有着无限的差异；它们的生命、行动和习性呈现出同样无法描述的多样性。因此，有些人试图为它们想象一些名字，其独特和奇异之处，如同烙印，标记出每种已知类的特征。有些，如鹰，被称为\OriginalTerm{裂翼类}{Schizoptera}，另一些，如蝙蝠，被称为\OriginalTerm{皮翼类}{Dermoptera}，再一些，如黄蜂，被称为\OriginalTerm{羽翼类}{Ptilota}，还有一些，如甲虫以及所有那些在壳和覆盖物中生出、打破监牢自由飞翔的昆虫，被称为\OriginalTerm{鞘翅类}{Coleoptera}。但我们有足够的日常用语来表征每个物种，并标记圣经在洁净与不洁净鸟之间所设立的区别。因此，食肉鸟类的种类是同一种类、同一构造，适合其生活方式：尖爪、弯喙、迅捷的翅膀，使它们容易俯冲捕食并在抓住后撕裂。那些捡拾种子的鸟的构造不同，而那些见什么吃什么的鸟的构造又不同。所有这些受造物中有多么大的多样性！有些是群居的，除了猛禽，它们除了配偶结合不知其他社会；但无数种类，如鸽子、鹤、椋鸟、寒鸦，喜欢共同生活。其中一些没有首领，生活在某种独立中；另一些，如鹤，不拒绝服从一个领袖。它们之间还有一个新的区别：有些是定居的、不迁徙的；另一些则做长途旅行，大部分在冬天临近时迁徙。几乎所有的鸟都可以被驯服并能接受训练，除了最弱小的，它们因恐惧和怯懦而不能承受手持续而烦人的接触。有些喜欢人类社会，居住在我们的住所；另一些则喜欢山野和旷野。它们特有的鸣叫声也有很大差异。有些唧唧喳喳，有些沉默，有些有悦耳嘹亮的声音，有些完全不和谐且不能歌唱；有些模仿人的声音，通过自然或训练学会模仿；另一些总是发出同样的单调叫声。公鸡骄傲；孔雀对自己的美貌虚荣；鸽子和家禽多情，总是寻找彼此的陪伴。鹧鸪奸诈嫉妒，为猎人捕捉猎物提供背信弃义的帮助。\par

4. 我曾说过，飞禽的行止与生活何等多样。其中一些无理性的受造物甚至拥有政府，倘若政府的特征在于使所有个体的活动都集中于一个共同的目的。这一点可在蜜蜂身上观察到。它们有共同的居所，在空中一齐飞翔，共同从事同样的工作；更奇异的是，它们在一王和监管者的引导下投入这些劳动，且不允许自己在未见国王领头飞翔的情况下飞向草地。至于这位国王，并非选举授予它权威——人民的无知常将最恶劣者置于权位；亦非命运——命运的盲目决断常将权柄交给最不配者；也不是世袭使牠登上王位——常见君王的子女因奢侈和奉承而堕落，对一切美德无知。乃是自然使牠成为蜂王，因为自然赋予了牠更大的体型、美丽和品性的温良。牠同其他蜜蜂一样有刺，但牠不用来报复。一条自然且不成文的法律原则是：被擢升至高位者，当在惩罚上宽厚。即使那些不效法其王的蜜蜂，也会立即为它们的鲁莽后悔，因为牠们随着刺而丧命。基督徒啊，你们听，你们被禁止\Quote{以恶报恶}，且被命令\Quote{以善胜恶}。以蜜蜂为榜样，牠建造蜂房时不伤害任何人，不干涉他人的财物。牠公开地用嘴从花朵上采集蜡，吸入散布于花上如露珠般的蜜，注入蜂房的空腔中。起初蜜是液态的，时间使之浓稠并赋予其甜味。\BibleBook{}{箴言}一书以称蜜蜂为智慧和勤奋的，给了牠最光荣和最高的赞美。牠为采集这宝贵的养料付出了多少辛劳，这养料使君王和庶民都得健康！在建造储藏蜂蜜的库房时，牠展现了何等巧艺和机巧！在将蜡铺成薄如薄膜的一层后，牠将其分布成相邻的隔室，这些隔室虽薄弱，却因数量和整体而支撑整个建筑。每个蜂房与其相邻者相连，由薄壁隔开；我们可见两三层蜂房彼此相叠。蜜蜂注意不造一个巨大的空腔，以免它在液态重量下破裂而泄漏。看，几何学的发现对聪明的蜜蜂来说不过是副产品！\par

蜂房的排状都是六边形，各边相等。它们不是直线互相支撑，以免支撑物挤压在空处而坍塌；而是下面六边形的角为上面蜂房的基础和基座，从而为下层质量提供稳固支撑，并使每个蜂房能安全地保存液态蜜。\par

5. 我们如何精确地检视鸟类生命的所有特点？夜间，鹤轮流守望；一些睡觉，另一些巡逻，为同伴谋得安静的睡眠。哨兵完成职责后鸣叫一声，便去睡觉；醒来的那一个则报答它所享有的安宁。同样的秩序也见于它们的飞行。一只在前领路，引领鹤群飞行一段时间后，便退至后方，将引领行军的职责交给后面的那只。\par

鹳的举止近乎理性的智慧。在这一地区，同一季节它们全部迁徙。它们以同一信号出发。在我看来，我们的乌鸦为它们担任护卫，前去迎回它们，并帮助它们抵御敌鸟的攻击。证据是，在这个季节不见一只乌鸦，它们带着伤痕回来，那是它们所提供帮助和援助的明显标志。是谁向它们解释了待客之道？是谁以离弃之罚威胁了它们？因为队伍中无一缺失。所有不接待客旅的人啊，你们关闭大门，冬季或夜间从不敞开房屋给行旅的人，请听。鹳对年迈者的关切，若我们的儿女反思，足以促使他们孝敬父母；因为没有一个人会愚昧到认为在德行上被无理性的禽鸟超越是可耻的事。鹳在父亲因年老羽毛脱落时环绕它，用翅膀温暖它，并充足地供养它；甚至在飞行中尽其所能帮助它，从两侧轻轻托起它，用翅膀扶持。这一行为如此著名，以至于感恩被称为\Quote{\OriginalTerm{报答慈恩之情}{antipelargosis}}。愿无人哀叹贫穷；愿家徒四壁者不因生计绝望，当他思虑燕子的勤劳。为筑巢，燕子用喙衔来草梗；因不能以爪衔泥，它先将翅尖在水中蘸湿，然后在极细的尘土中翻滚，从而获得泥。它以此泥如胶般将草梗逐一黏合后，喂养幼雏；若有雏鸟眼睛受伤，它有天然的药方治愈幼鸟的视力。\par

这景象当警告你，莫因贫穷而走入恶途；即使你陷入绝境，也不可失去一切希望；不可自弃于怠惰与游手好闲，而要投靠天主。若祂对燕子如此慷慨，那对于那些全心呼求祂的人，祂岂不更会如何？\par

翠鸟是一种海鸟，它在岸边产卵，或将蛋产于沙中。它在仲冬产卵，此时暴风将海浪拍向陆地。然而，在翠鸟孵卵的七天里，一切风都平息了，海浪也平静了。\par

因为孵化幼鸟只需七天。它们需要食物才能成长，天主出于祂的慷慨，又赐予这微小的动物七天。所有航海者都知道这一点，并称这些日子为太平日子。既然天主眷顾为无理性的受造物设立了这些奇妙的律法，是为了促使你向天主祈求你的救恩。祂岂不愿为你——那按祂肖像受造的你——行任何奇事？当为了这么一只小鸟，那浩瀚可怕的大海被约束，在严冬中被命令平静下来。\par

6. 据说斑鸠一旦与伴侣分离，就不会缔结新的结合，而是为纪念最初的盟约保持寡居。妇女们，请听！即使在无理性的受造物中，对寡居何等尊重，它们竟宁愿寡居也不愿有不体面的多婚。老鹰在抚育幼雏时表现出最大的不公。当它孵出两只雏鹰时，它会用翅膀击打并推出一只丢在地上，只承认剩下的那只。是觅食的困难使它抛弃自己生出的后代。但据说鱼鹰不会让它死去，而是将其带走与自己的幼雏一同抚养。这就是那些以贫穷为借口遗弃子女的父母，也是在分配遗产时作出不平等划分的父母。既然他们给予每个子女同样的存在，理应以同等而不偏袒的方式供给他们生计。当心不要模仿钩爪禽鸟的残忍。当它们看到幼雏已能迎空飞翔时，便将它们推出巢外，用翅膀击打推挤，不再给予丝毫照顾。乌鸦对幼雏的爱值得赞美！当它们开始飞行时，乌鸦跟随它们，喂食它们，并长时间提供营养。许多鸟无需与雄性结合就能受孕，但它们的蛋是不结果的，除了秃鹫的蛋——据说秃鹫更常在不交配的情况下生育，尽管它们寿命很长，常达百年。我请求你们注意并记住鸟类史中的这一点：如果你们曾见有人嘲笑我们的奥迹，认为童贞女成为母亲而不失童贞的纯洁是不可能且违反自然的，请想想那愿意借着愚妄的宣讲拯救信靠者\BibleRef{格前 1:21}，已事先在自然中给了我们上千个相信这奇事的理由。\par

7. \Quote{\Emph{水要滋生有生命的爬行动物，并有飞鸟飞在地面以上，天空的穹苍之中。}}它们奉命飞行在地面以上，因为大地供给它们食物。\Quote{天空的穹苍之中}，即如前所述，在空气的那部分，称之为\OriginalTerm{天}{\GreekText{οὐρανός}}，由\OriginalTerm{视}{\GreekText{ὁ} \GreekText{ρᾶν}}一词而来；称为穹苍，因为在我们头顶延伸的空气，与以太相比，密度更大，并由从大地升起的雾气所增厚。于是你有了装饰的天、美丽的地、充满自身受造物的海、以及充满四处飞翔飞鸟的空气。好学的听者，请想想天主从无中造出的这些受造物，请想想我的讲论为避免冗长及不超出限度而省略的一切；请处处认出天主的智慧；永不止息地惊叹，并借着每一个受造物荣耀造物主。\par

有些鸟类夜间生活在黑暗中；另一些则白天在充足光线下飞翔。蝙蝠、猫头鹰、夜鸦是夜鸟：如果你偶然无法入睡，请反思这些夜鸟及其特性，并荣耀它们的造物主。夜莺在孵卵时何以总是醒着，整夜不断歌唱？蝙蝠这种动物何以同时是四足动物和飞禽？何以它是鸟类中唯一有牙齿的？何以它像四足动物一样胎生，却能腾空飞起，不是靠翅膀，而是靠一种薄膜？蝙蝠彼此之间有何等自然的爱！它们如何像链条一样交织，一个悬挂在另一个之上！这在人类中是极其罕见的景象，因为人类大多更喜欢个人和私下的生活，而非共同生活的结合。那些献身于虚妄科学的人岂不有猫头鹰的眼睛？猫头鹰在夜间视力敏锐，却被太阳的光辉眩目；同样，这些人的理智虽善于凝思虚妄，却在真光面前盲目。\par

白天，你随处也容易赞叹造物主！看，家鸡以尖声啼叫唤你工作，它如同太阳的先驱，又像早起的旅人，催促农夫去收割！鹅何等警觉！它们以何等敏锐的嗅觉察觉隐藏的危险！它们不是曾拯救过那座帝国城市\Place{罗马}吗？当敌人通过地下通道潜入以夺取罗马的议会大厦时，鹅不是发出了警报吗？有哪一种鸟的本性不令我们赞叹？谁向秃鹫宣告，当人们列阵相战时将有杀戮？你可以看到成群的秃鹫跟随军队，推算战事准备的结果——这是一种非常接近于人类理性的推算。我怎能向你描述蝗虫可怕的入侵？它们一旦得到信号就从各处升起，在全国各地扎营。它们不会攻击庄稼，除非领受了天主的命令。或者我该描述这灾祸的解救者——鸫鸟，它以永不满足的食欲追随着蝗虫，以及慈爱的天主为怜悯人类所赐予它的吞噬本性？蚱蜢如何调节其歌声？为何它在正午时分因吸气而扩张胸部时更加悦耳？\par

但在我看来，当我想要描述飞禽的奇迹时，我落后得比我的双脚试图追赶它们飞翔的快速还要甚远。当你看到蜜蜂、黄蜂，简而言之，所有那些被称为昆虫的飞行生物，因为它们全身有一道切割线，请思考它们既没有呼吸也没有肺，并且它们通过全身各部分被空气支持。因此，如果它们被油覆盖，就会死亡，因为油堵塞了它们的毛孔。用醋冲洗它们，毛孔重新打开，动物复活。我们的天主没有创造任何多余的东西，也没有省略任何必要的东西。如果你现在注视水中的生物，你会发现它们的构造完全不同。它们的脚不像乌鸦那样分叉，也不像肉食动物那样带钩，而是大而呈膜状；因此，它们能轻易游泳，用脚上的膜像桨一样推水。注意天鹅如何将脖颈潜入水的深处以从中觅食，你就会理解造物主的智慧：祂给这种生物一个比脚还长的脖子，使它能够像抛线一样投出脖子，并取走藏在深水底的食物。\par

8. 如果我们只读圣经的话，我们只找到几个简短的音节。\BibleQuote{让水生出飞鸟，使它们飞行在地面之上，在天空的穹苍之中}，但如果我们探究这些话语的意义，那么造物主智慧的伟大奇迹便显现出来。祂在飞禽之间预见了何等大的差异！祂如何按种类划分它们！祂如何用独特的品质来表征每一种！然而，一整天都不足以让我讲述空中的奇迹。大地正在召唤我去描述野兽、爬行动物和牲畜，它准备向我们展示与植物、鱼和鸟相媲美的景象。\BibleQuote{让大地生出活魂}，即家畜、野兽和爬行动物，各按其类。你们这些不相信保禄在复活中所应许的改变的人，当你们看到空中生物中有如此多的变形时，还有什么可说的呢？关于印度的带角蠕虫，我们被告知了什么！它先变为毛毛虫，然后变成嗡嗡叫的昆虫，并不满足于这种形态，它用松散宽阔的甲片而非翅膀来包裹自己。因此，女人们，当你们坐着忙于纺织时，我是指由中国送来给你们做精致衣服的丝绸，请记住这个生物的变形，对复活形成一个清晰的概念，并且不要拒绝相信保禄宣告的为所有人的改变。\par

但我惭愧地看到我的讲论超出了通常的限度；若我考虑刚才向你们讲述的丰富内容，我感到自己正被带出界限；但当我反思创造之工中显示的无尽智慧时，我似乎只是在我故事的开端。然而，我并没有毫无益处地久留你们。因为你们直到晚上能做什么呢？你们没有客人催促，也没有宴会在等。那么，就让我利用这身体上的斋戒来使你们的灵魂喜乐。你们常常为享乐而服侍肉身，今天要持久地服侍灵魂。\BibleQuote{你要在上主内喜乐，祂必赐给你心之所欲。} 你们喜爱财富吗？这里有属灵的财富。\BibleQuote{上主的典章是真实公义的，全都正直。比金子可羡慕，比精金可羡慕。} 你们喜爱享乐和欢娱吗？请看上主的圣言，对于健康的灵魂，它们\BibleQuote{比蜂蜜和蜂房更甘甜。} 如果我让你们离去，如果我解散这会众，有些人就会跑去掷骰子，在那里他们会遇到污言秽语、悲伤的争吵和吝啬的痛苦。魔鬼站在那里，用有斑点的骨头点燃赌徒的狂怒，将同一笔钱从桌子的一边移到另一边，一会儿用胜利使一人得意，使另一人陷入绝望，一会儿使第一个充满夸耀，使他的对手狼狈不堪。身体上的斋戒与使灵魂充满无数邪恶有何益处？不赌博的人会在别处消磨闲暇。从他的口中吐出什么轻浮之言！他的耳朵听到什么蠢话！没有敬畏上主的闲暇，对于那些不知道时间价值的人，是邪恶的学校。我希望我的话有益处；至少通过在这里占用你们的时间，阻止了你们犯罪。因此，我留你们越长，你们就越远离邪恶之路。\par

一个公正的审判者会认为我已经说得足够，不是看他是否考虑创造之富，而是看他想及我们的软弱以及在快乐之事上应持守的限度。大地以其植物欢迎你，水以其鱼欢迎你，空气以其鸟欢迎你；大陆也正准备好向你提供同样丰盛的宝藏。但让我们结束这清晨的盛宴，免得饱足削弱你对晚间盛宴的胃口。愿祂以祂创造之工充满万有，并在各处留下祂奇迹的可见纪念，以在耶稣基督、我们的主内的所有属灵喜乐充满你们的心，愿光荣和权能归于祂，永无穷尽。阿们。\par

\SourceRecord{Translated by Blomfield Jackson.；Nicene and Post-Nicene Fathers, Second Series；Vol. 8.；Edited by Philip Schaff and Henry Wace.；Buffalo, NY: Christian Literature Publishing Co.,；1895.；<http://www.newadvent.org/fathers/32018.htm>.}

% structure: p0001 source=h1 target=section reason=page-title

\section{六天创世（第九篇讲道）}

\Emph{地上动物的创造。}\par

1. 你们觉得我早晨讲道的滋味如何？在我看来，我就像一个办筵席的穷人，虽有好意，却因缺乏昂贵的菜肴而使宾客扫兴。他以简陋的菜肴慷慨地摆满桌子，但他的野心只显露出愚拙。你们来判断我是否遭遇了同样的命运。然而，无论我的讲道如何，你们要小心，不可轻视它。没有人拒绝坐在\Person{厄里叟}的席上，而他只给朋友们野生的蔬菜。\BibleRef{列下 4:39}我懂得寓意的法则，虽然更多是从他人的著作中得来，而非靠我自己。的确有些人，不接受圣经的字面意义，对他们来说，水不是水，而是另一种本性；他们在植物、鱼身上看到他们幻想所愿的东西；他们改变爬虫和野兽的本性以适应他们的寓意，就像解梦者解释睡梦中的异象以服务于自己的目的。对我来说，草就是草；植物、鱼、野兽、家畜，我都按字面意义理解。\Quote{\BibleQuote{我不以福音为耻。}}\BibleRef{罗 1:16}那些论述宇宙本性的人已详细讨论了地球的形状。它是球形的还是圆柱形的，像圆盘一样处处均匀地圆，还是像扬谷的簸箕，中间凹陷？宇宙学家提出了所有这些猜测，每个人都推翻其前人的理论。天主的仆人\Person{梅瑟}对形状保持沉默，这并不会使我对宇宙的创造减少重视；他没有说地球的周长是十八万弗隆；他没有测量太阳绕地球运行时它的影子投射到空气中多远，也没有说明这个影子投射到月亮上如何产生月食。他把所有对我们无关紧要的东西都当作无用的而略过了。那么，我岂应把愚拙的智慧置于圣神的谕言之上？我岂不应颂扬祂，祂不愿以这些虚妄充满我们的心，而为了我们灵魂的建立和成全安排了圣经的全部计划？在我看来，那些沉迷于寓意扭曲意义的人没有理解这一点，他们试图以自己杜撰的威严来装点圣经。这是相信自己比圣神更智慧，并以释经为借口提出自己的见解。让我们按圣经所写的来听吧。\par

2. \Quote{\Emph{让大地生出活物。}}\BibleRef{创 1:24}看，天主的话渗透了受造界，从那时起就开始运行那今日显见的效能，并将运行到世界的终结！就像一个被推的球，如果遇到下坡，就会凭借它的形状和地面的性质滚下去，直到到达平坦的地面才停止；同样，本性一旦被天主的命令推动，就以同等的步伐穿越受造界，通过诞生和死亡，保持种类的相似性，直到终末。本性总是让马生马，狮生狮，鹰生鹰，并藉着这些不间断的传承保全每个动物，将其传递到万物的终结。动物不会因时间的流逝而看到其特性被破坏或抹去；它们的本性，如同刚刚被造一样，随岁月的流逝而保持常新。\Quote{让大地生出活物。}这道命令一直延续，大地也没有停止服从创造者。因为，如果有动物是依次由它们的先辈产生的，那么也有其他动物我们今天看到直接从大地生出。在潮湿天气里，大地生出蚱蜢和无数飞在空中的昆虫，它们因太小而没有名字；她也生出老鼠和青蛙。在\Place{底比斯}（埃及）附近，天气炎热时在大雨之后，田野上布满了田鼠。我们看到仅淤泥就产生鳗鱼；它们并非来自卵，也不来自其他方式，而是大地独自生出它们。让大地产出活物。\par

地上的牲畜，头朝下向着大地。人，属天的生长物，无论在身体构型还是灵魂尊严上都远超它们。四足动物的形态如何？它们的头俯向大地，注视着肚腹，只追求肚腹的益处。人啊！你的头朝向天空；你的眼睛向上看。因此，当你因肉体情欲而贬低自己，作肚腹和低贱部分的奴隶时，你就接近无理性的禽兽，变得像它们一样。你被召叫追求更高贵的事业：\Quote{\BibleQuote{你们要思念上面的事，那里有基督坐在天主右边。}}\BibleRef{哥 3:1}使你的灵魂超出尘世；从你身体的自然构造中汲取行为准则；把你的思念安放在天上。你真正的家乡是天上耶路撒冷；你的同胞和同乡是\Quote{\BibleQuote{那些被记录在天上的首生者。}}\BibleRef{希 12:23}\par

3. \Quote{\Emph{让大地生出活物}}。因此，当野兽的灵魂出现时，它并非隐藏在地内，而是因天主的命令而生。野兽拥有同一类的灵魂，其共同特征是无理性。但每种动物都以特有的品质相区别。牛稳重，驴懒惰，马情欲强烈，狼不可驯服，狐狸诡诈，鹿胆怯，蚂蚁勤劳，狗在友谊中感恩而忠诚。每种动物受造时，其本性的独特特征便按适当比例显现在它身上：狮子有精神、喜好独居、不合群的性格。他是动物中真正的暴君，出于天生的傲慢，只允许少数动物分享他的荣耀。他鄙视昨日的食物，从不返回剩余的猎物。大自然赋予他的发声器官如此巨大的力量，以至于常常有速度更快的动物仅因他的吼叫就被捕获。豹子跳跃猛烈而冲动，其身体适合其活跃和轻盈，与灵魂的运动相协调。熊有迟钝的本性、独特的方式、狡猾的性格且极为隐秘；因此它有相应的身体——沉重、厚实、没有关节，正如冰冷的穴居者所需。\par

当我们思虑这些无理性受造物对自己生命所具有的天生与固有的关怀时，我们将会受到激励去照管自己，并思考我们灵魂的救恩；或者更确切地说，当我们被发现甚至不如模仿野兽时，我们将更被定罪。熊常常受重伤，却能照料自己，巧妙地用毛蕊花（一种性质极为收敛的植物）填满伤口。你也会看到狐狸用松树的落果治愈伤口；乌龟饱食蝰蛇的肉后，借助墨角兰的功效来对抗这种有毒动物；蛇通过吃茴香治愈眼疾。\par

在应对气候变化方面，理性智慧岂不也被动物所超越吗？我们岂不见到羊在冬天临近时贪婪地吃草，仿佛为未来的匮乏做准备？我们岂不也见到牛在冬季长期关在栏中，通过自然感觉认出春天的回归，并一致地把头转向马厩尽头的门口？细心的观察者曾注意到刺猬在其洞穴的两端各开一个口。如果北风将要吹来，它就堵住朝北的洞口；如果南风随后吹来，它就转到北边的门口。这些动物教训人什么？它们不仅向我们显示造物主对万物的眷顾，也显示野兽甚至对未来有某种预感。因此，我们不应执着于现世生命，而应密切关注来世。人啊，你难道不为自己勤劳？你难道不效法蚂蚁的榜样，在此世为来世储蓄安息？蚂蚁在夏季为冬天积蓄宝物。它绝不纵容自己懒惰，在严冬感到它的严厉之前，它就以不可战胜的热忱急忙工作，直到充足地装满它的仓库。在此，它又是何等远离疏忽！它以何等明智的远见来管理，以使其储存尽可能持久！它用钳子将谷粒切成两半，以防它们发芽而不再能作为食物。如果谷粒潮湿，它就晒干；它并非在任何天气下都摊开谷粒，而是在感到空气将保持温和时。你确可肯定，只要蚂蚁将谷粒摊在外面，你就绝不会看到雨水从云中降下。\par

什么言语能领会造物主的奇妙？什么耳朵能理解它们？又有什么时间足以述说它们？那么，让我们与先知一同说：\BibleQuote{上主，你的化工何其繁多！你以智慧创造了它们全体}。我们将无法在自辩时说，我们从书本中学到了有用的知识，因为大自然未经教导的律法使我们选择有利于我们的事物。你知道你应当向邻人行什么善吗？就是你自己期望从他那里得到的善。你知道什么是恶吗？就是你不愿他人对你做的恶。植物学研究和药草经验并未使动物发现那些对它们有用的东西；而是每个动物自然地知道什么是有益的，并奇妙地取用适合其本性之物。\par

4. 美德也自然地存在于我们内，灵魂与它们有亲和性，并非通过教育，而是通过大自然本身。我们不需要课本来憎恨疾病，而是我们自行驱除折磨我们的事物；灵魂不需要师傅来教导我们避免邪恶。如今，一切邪恶都是灵魂的疾病，正如美德是灵魂的健康。因此，那些把健康定义为自然功能有序运作的人，他们的定义是恰当的；这一定义可以无所畏惧地应用于灵魂的良好状态。因此，灵魂无需课本书，便能自行达到适合并合乎本性的境界。由此，节制处处受赞扬，正义受尊敬，勇气受钦佩，明智为众目的所向；这些美德关乎灵魂，胜过健康关乎身体。孩子们，要孝敬父母；你们作父母的，\BibleQuote{不要惹你们的子女发怒}。大自然不也这样说吗？保禄教导我们并非新事；他只是巩固了自然的联系。如果母狮爱其幼崽，如果母狼为保护幼崽而战斗，那么违背诫命、违反自然本身的人——或是辱骂父亲年老的儿子，或是因再婚而忘记头生子女的父亲——又该说什么呢？\par

在动物中，不可战胜的情感使父母与子女结合在一起。是天主，造物主自己，以感情的力量取代了它们中的理性。由此，一只羔羊从羊圈跳出，在千只羊中认出它母亲的颜色和声音，跑向她，寻求自己的奶源。如果它母亲的乳房干涸，它便满足，而不停步，经过更丰盛的乳房。母亲又如何从众多羔羊中认出它？所有的声音相同，颜色相同，气味相同，至少就我们的嗅觉而言。然而，这些动物有一种比我们的感知更微妙的感官，使它们认出自己的。小狗还没有牙齿，但仍用嘴自卫，抵抗任何逗弄它的人。小牛还没有角，但它已经知道它的武器将在哪里生长。这里我们有明显的证据，证明动物的本能是天生的，并且所有受造物中没有任何无序或不可预见之事。一切都带着造物主智慧的印记，并表明它们来到生命时已具备了确保自身保存的手段。\par

狗并未分得一份理性；但在它身上，本能具有理性的力量。狗从自然中学会了复杂推理的奥秘，世上的智者经过多年的研究也几乎无法解开。当狗追踪猎物时，如果他看到猎物分向不同的方向，他便检查这些不同的路径，唯有言语未能宣布他的推理。这个生物，他说，去了这里或那里或另一个方向。既不是这里也不是那里；因此是第三个方向。这样，忽略了错误的踪迹，他发现了真正的踪迹。那些认真忙于论证理论，在尘土上画线，并拒绝两个命题以表明第三个为真的人，所做的又有何更多呢？\par

狗的感激难道不使所有忘恩负义于恩人的人羞愧吗？据说许多狗在荒僻之处为主人被谋杀而殉死。其他的狗，当罪行刚刚发生时，曾带领搜索凶手的人，并使罪犯被绳之以法。那些不仅不爱创造并养育他们的主，反而与用口攻击主的人为友，与他们同席，并在分享食物时亵渎赐予食物者的人，将说什么呢？\par

5. 但让我们回到创造的景象。最容易捕捉的动物是最多产的。正因如此，野兔和野山羊生许多幼崽，野绵羊生双胞胎，唯恐这些物种被食肉动物消耗而消失。相反，猛兽只产少数，母狮艰难地产下一只狮子；因为，如果他们说真的，幼狮是用爪子撕裂母亲而出的；而毒蛇只有咬穿子宫才出生，给母亲施加适当的惩罚。因此，自然界中一切都被预见，一切都在持续的照料之下。如果你甚至检查动物的肢体，你会发现造物主没有给它们任何多余的东西，也没有遗漏任何必要的东西。给食肉动物尖利的牙齿，它们的本性需要这些来维持生计。那些牙齿配备不全的动物有几个不同的食物储存器官。由于食物在第一个器官中未被充分分解，它们具有在吞咽后将食物返回的能力，直到被反刍碾碎后才被吸收。反刍动物的第一、第二、第三和第四胃并非闲置；它们每一个都履行必要的功能。骆驼的脖子长，以便它能低头到脚够到它赖以生存的草。熊、狮、虎，所有这类动物，脖子短，埋在肩膀里；这是因为它们不靠草为生，无需弯向地面；它们是食肉动物，吃它们捕食的动物。\par

大象为何有鼻子？这巨兽乃陆地动物中最庞大者，受造是为令遇见者惊惧，其体自然硕大而多肉。若其颈粗大并与腿相称，则难驾驭，且过重使其倾向地面。事实上，其头以短椎骨连接脊背，以鼻代颈，以此取食汲饮。其腿无关节，如并立之柱，支撑体重。若腿松软有节，则关节常会弯折，无论欲跪或立，均无法支撑体重。然其足下有一小踝关节，代替腿与膝之关节，若后者活动则无法承受这巨大摇摆之躯。故需此几及足之鼻。\newline{}汝可曾见它们在战争中如活塔般行进于方阵前列，或以不可抵挡之冲锋如肉山般冲破敌军阵营？若其下部不与体型相称，则绝不能支撑自身。\newline{}据说象活三百余年，此乃其有坚实无节之足的另一原因。然如我们所说，其鼻形似蛇且灵活，从地取食并举起。故此我们有理由说，受造界中无任何多余或欠缺之物。\newline{}天主竟将此怪兽驯服于我们，甚至能领会教训、忍受我们之击打，此乃明显证据，证明造物主因我们系按祂肖像受造，已将万有置于我们治理之下。\newline{}不仅在大动物身上可见无与伦比之智慧，在最小者中亦见同样奇事。高山之巅近云常受风击、终年严冬者，并不比深谷更能激发我赞叹，因深谷避过高峰风暴而保持恒久温和。同样，在动物构造中，我惊叹大象之大，亦不亚于鼠（象所畏者）或蝎之细刺——至高工匠使之如管道中空以注入毒液。\newline{}莫让人指责造物主创造毒物、害人敌人。否则让他们视学监为有罪，因其用杖鞭管教少年躁动以维持秩序。\par

6. 野兽为信德作证。你信赖主吗？\Quote{你要践踏蝮蛇和蜥蜴，将狮子和毒龙踏在脚下。} 藉着信德，你有能力践踏蛇和蝎子。你岂不见那蝮蛇缠住保禄的手，当他拾柴时，并未伤害他，因为发现这位圣人充满信德吗？若你没有信德，与其惧怕野兽，不如惧怕你自身的不信，它使你易受一切败坏。但我看出，长久以来，你们一直问我关于人的受造的解说，我想我可以听见你们心中呼喊：我们正在受教于我们所有物的本性，却不认识我们自己。那么，让我来说吧，因为这是必要的，让我结束我的犹豫。诚然，最难的科学是认识自我。不仅我们的眼睛——它不能见到自身，尽管外界无一能逃过它；而且我们的心智，它尖锐地发现他人的罪过，却迟迟不能认出自己的过失。因此，我的言论，在热切探求我自身以外的事物之后，在探索我自己的本性时却迟缓而犹豫。然而，观察天地并不比细察我们的存在使我们更认识天主；先知说：\Quote{我受造，令人惊惧，奇妙莫测；} 也就是说，在观察自己时，我认识了祢的无限智慧。天主说：\BibleQuote{让我们造人。} \BibleRef{创 1:26} 神学的光芒岂不像透过窗户一样，藉这些言语照耀？第二位岂不奥秘地显示出自己，却尚未在大日之前显明？那抗拒真理、佯称天主向自己说话的犹太人在哪里？据说是祂说了话，又是祂创造了。\Quote{要有光，就有了光。} 但这样他们的话包含明显的荒谬。铁匠、木匠、鞋匠，若无人帮助，独自在他的工具前，岂会对自已说：让我们造剑，让我们装配犁，让我们做靴子？他岂不沉默地从事他的手艺？何等荒谬的蠢行，说有人坐下来命令自己、监视自己、约束自己、催促自己，用主人的口气！但不幸的人们竟敢毁谤主自己。他们那惯于说谎的舌头还有什么不敢说呢？然而，在此，话语堵住了他们的口：\BibleQuote{让我们造人。} 请告诉我：那么只有一个位格吗？经上并未说：\Quote{让人受造，} 而是说：\BibleQuote{让我们造人。} 神学的宣讲在受训者出现之前仍笼罩在阴影中，但如今，人的创造即将来临，信德显现自身，真理的信条在全部光明中出现。\BibleQuote{让我们造人。} 基督的仇敌啊，请听天主对祂的同工说话，对那藉着祂也创造了诸世界、以祂权能的话支撑万有者说话。但祂并未让真宗教的声音无回应。因此，犹太人，这敌视真理的民族，当发现自己受逼迫时，如同被激怒的野兽攻击人，在笼子的栅栏旁咆哮，显出其本性的残忍与凶暴，却无法平息其怒火。他们说：天主是向多位说话；祂是对祂面前的众天使说：\BibleQuote{让我们造人。} 犹太人的虚构！一则寓言，其轻浮显出它的来源。为否定一位，他们承认众多。为否定圣子，他们提升仆人到谋士的尊严；他们使我们的同仆成为我们受造的代理人。完美的人达到天使的尊严；但受造物岂能与造物主相似？请听下文。\BibleQuote{按我们的肖像。} 你如何回答？天主与天使有同一肖像吗？圣父和圣子因绝对必然性而有同一形式，但这形式在此是按神圣的方式理解，不是身体的形状，而是天主性的特有属性。请听，你们这些属于新割损派的人\BibleRef{斐 3:2}，你们在基督信仰的外表下，坚固了犹太人的错误。祂说：\BibleQuote{按我们的肖像，} 是对谁说的，如果不是对那一位，祂是\BibleQuote{祂荣耀的光辉和祂本体的真像}\BibleRef{希 1:3}，\BibleQuote{不可见天主的肖像}\BibleRef{哥 1:15}？那么，是对祂的活的肖像，对那曾说\BibleQuote{我与父原是一体}\BibleRef{若 10:30}、\BibleQuote{看见了我就是看见了父}\BibleRef{若 14:9}的那一位，天主说：\BibleQuote{让我们按我们的肖像造人。} 在这些仅有同一肖像的存在者中，哪里有不相似？\BibleQuote{于是天主造了人。}\BibleRef{创 1:27}经上不是说：\Quote{他们造了。} 在此，圣经避免位格的复数。在启发了犹太人之后，它消除了外邦人的错误，将自己置于合一的庇护之下，使你明白圣子与圣父同在，并保护你免于多神论的危险。祂照天主的肖像造了他。天主仍向我们展示祂的同工，因为祂不说：\BibleQuote{照祂的肖像，} 而是：\BibleQuote{照天主的肖像。}\par

若天主允许，我们稍后将说人是如何按天主的肖像受造的，以及他是如何分享这相似。今天我们只说一句话。如果只有一肖像，那么那不堪忍受的亵渎——妄言圣子不像圣父——从何而来？何等的忘恩负义！你自己接受了这相似，却拒绝赐予你的恩主！你佯称个人保住那在你内是恩宠恩赐之物，却不愿圣子保持祂与生祂者之间的自然相似。\par

但是，那早已把太阳送到西边的黄昏，却迫使我缄默。因此，就让我满足于已说的话，让我的讲论安歇吧。到此为止，我已向你们讲得足够，足以激发你们的热情；借圣神的助佑，我将为你们更深入地探究随后要讲的真理。那么，我恳请你们，喜爱基督的会众，满怀喜乐地退去吧，并以对我话语的怀念来装饰和圣化你们的餐桌，而不是用各种佳肴美味。愿\OriginalTerm{阿诺莫伊派}{Anomœan}蒙羞，愿犹太人受辱，愿信友因真理的教义而欢欣，愿主受光荣，愿荣耀与权能归于主，至于无穷世。阿们。\par

\SourceRecord{Translated by Blomfield Jackson.；Nicene and Post-Nicene Fathers, Second Series；Vol. 8.；Edited by Philip Schaff and Henry Wace.；Buffalo, NY: Christian Literature Publishing Co.,；1895.；<http://www.newadvent.org/fathers/32019.htm>.}
